\documentclass{article}

\usepackage {tikz}
\usepackage[utf8]{inputenc}
\usepackage[mono=false]{libertinus}
\usepackage[T1]{fontenc}
\usepackage{sectsty}
\allsectionsfont{\sffamily}
\usepackage{tocloft}
\usepackage{multicol}

\renewcommand{\cftsecfont}{\bfseries\sffamily}
\usepackage[numbers]{natbib}
\usepackage{enumitem,lipsum}
\usepackage{graphicx}
\usepackage{subcaption}
\usepackage{geometry}
\usepackage{mathtools}
% \geometry{margin=1in}
\usepackage{amsmath}
\usepackage{amssymb}
\usepackage{amsthm}
\usepackage{fancyhdr}
\usepackage{array}
\usepackage{hyperref}
\usepackage[breakable]{tcolorbox}
\usepackage{bbm}
\usepackage{parskip}
\usepackage{physics}
\usepackage[hang,flushmargin]{footmisc}
\usepackage{mathtools}
\usepackage{stmaryrd}
\DeclareSymbolFont{rsfso}{U}{rsfso}{m}{n}
\DeclareSymbolFontAlphabet{\mathscr}{rsfso}
\usepackage{tikz}
\usepackage{pgfplots}
\pgfplotsset{compat=1.15}
% \usepackage{mathrsfs}
\usetikzlibrary{arrows}
\usepackage{comment}
\usepackage{blkarray}
% \usepackage{ntheorem}
\usepackage{multirow}
\usepackage{amsmath,amssymb,graphicx,amscd}
\usepackage{mathrsfs}
% pour les S des spheres
\usepackage{mathrsfs}
\usepackage{dsfont}
\usepackage{stmaryrd}
\usepackage{wrapfig}
\usepackage{xcolor}
\usepackage{pstricks}
\usepackage{lastpage}
\hypersetup{
    colorlinks=true,
    linkcolor=magenta,
    filecolor=magenta,      
    urlcolor=cyan,
    citecolor = magenta
}

\urlstyle{same}

\newtheorem{thm}{Theorem}[section]
\newtheorem{lem}[thm]{Lemma}
\newtheorem{prop}[thm]{Proposition}
\newtheorem{ass}{Assumption}

\theoremstyle{definition}
\newtheorem{definition}{Definition}[section]

\theoremstyle{remark}
\newtheorem*{remark}{Remark}
\newtheorem*{hint}{Hint}
\newtheorem*{note}{Note}

\def\Xint#1{\mathchoice
{\XXint\displaystyle\textstyle{#1}}%
{\XXint\textstyle\scriptstyle{#1}}%
{\XXint\scriptstyle\scriptscriptstyle{#1}}%
{\XXint\scriptscriptstyle\scriptscriptstyle{#1}}%
\!\int}
\def\XXint#1#2#3{{\setbox0=\hbox{$#1{#2#3}{\int}$ }
\vcenter{\hbox{$#2#3$ }}\kern-.6\wd0}}
\def\ddashint{\Xint=}
\def\dashint{\Xint-}

% declare a new theorem style
\newtheoremstyle{mystyle}%
{3pt}% Space above
{3pt}% Space below 
{\itshape\color{red}}% Body font
{}% Indent amount
{\bfseries\color{red}}% Theorem head font
{.}% Punctuation after theorem head
{.5em}% Space after theorem head
{}% Theorem head spec (can be left empty, meaning ‘normal’)
\theoremstyle{mystyle}
\newtheorem*{grade}{Comment}
\newcommand{\borel}{\mathcal{B}}
\newcommand{\expo}{\mathrm{Exp}\,}
\newcommand{\ber}{\mathrm{Ber}\,}
\newcommand{\binomial}{\mathrm{Bin}\,}
\newcommand{\negbinomial}{\mathrm{NB}\,}
\newcommand{\Var}{\mathrm{Var}} %variance
\newcommand{\Geom}{\mathrm{Geom}\,}
\newcommand{\Pois}{\mathrm{Pois}\,}
\newcommand{\Erlang}{\mathrm{Erlang}\,}
\newcommand{\Det}{\mathrm{det}}
\newcommand{\sgn}{\mathrm{sgn}}
\newcommand{\Cov}{\mathrm{Cov}}
\renewcommand{\Pr}{\mathbb{P}}
\newcommand{\E}{\mathbb{E}}
\newcommand{\KL}{\mathrm{KL}\,}
\newcommand{\indep}{\perp \!\!\!\! \perp}
\newcommand{\1}{\mathbbm{1}}
\newcommand{\CC}{{\mathbb C}}
\newcommand{\HH}{{\mathbb H}}
\newcommand{\R}{\mathbb R}
\newcommand{\QQ}{{\mathbb{Q}}}
\newcommand{\NN}{{\mathbb{N}}}
\newcommand{\ZZ}{{\mathbb{Z}}}
\renewcommand{\P}{{\mathbb{P}}}
\newcommand{\rd}{{\mathrm{d}}}
\newcommand{\Prob}{{\mathbb{P}}}

 \newcommand{\I}{1\!\!1}                   %Indicator function

 \newcommand{\ii}{{\mathrm{i}}}
  \newcommand{\Id}{{\mathrm{Id}}}
\newcommand{\law}{\overset{\mbox{\rm \scriptsize law}}{=}}
\newcommand{\convlaw}{\overset{\mbox{\rm \scriptsize law}}{\longrightarrow}}
\let\oldproof\proof
\renewcommand{\proof}{\color{blue}\oldproof}

% \renewcommand{\chaptermark}[1]{\markboth{#1}{#1}}
% \renewcommand{\sectionmark}[1]{\markright{#1}}
\pagestyle{fancy}
\fancyhf{}
\renewcommand{\headrulewidth}{0pt}
\fancyhead[C]{}
\fancyhead[L]{\textit{\leftmark}}
% \fancyhead[RE]{\textit{\leftmark}}
\fancyhead[R]{ \quad \thepage of \pageref{LastPage}}
\fancypagestyle{plain}{
    \fancyhead{}
    \renewcommand{\headrulewidth}{0pt}
    \fancyfoot[C]{\textbf{\thepage}}
}

%-----------------------------------------------------------------

\title{Diffusion policy}
\author{Zhen Bao}

\begin{document}

\maketitle
% \section{Notations\& Definitions \& previous results}
% Define $J(\theta)\coloneqq V^{\pi_\theta}(\rho)=\E^\pi_\rho \left[\sum^\infty_{h=0}\gamma^hr(x_h,a_h)\right]$, where $a_h\sim \pi_\theta(\cdot|x_h)$ and $x_{h+1}\sim \P(\cdot|x_h,a_h)$. \\
% Define $\P^{\theta,x}$ s.t. $\P^{\theta,x}\circ(\bar{a}_T)^{-1}=\P\circ(\bar{a}_T^x)^{-1}$ clear definition by Girsanov. This only for reference.\\
% Policy gradient $\nabla_\theta J(\theta)=\frac{1}{1-\gamma}\int_\mathcal{X}\E\left[S_\theta(x)Q^{\pi_\theta}(x,\bar{a}^x_t)\right]d^{\pi_\theta}_\rho(dx)$\\
% $G_H(f)=\sum^{\infty}_{j=0}\gamma^{j}r(X^f_j,A^f_j)$; Truncated: $G_H(f)=\sum^{H-1}_{j=0}\gamma^{j}r(X^f_j,A^f_j)$. With $f$ parametrization with $\theta$, let $G_H(\theta,\xi)=\sum^{H-1}_{j=0}\gamma^{j}r(X^{f_\theta}_j,A^{f_\theta}_j)$. $PG(\theta)=\E_\xi[G(\theta,\xi)]$

% \section{Proof structure}
% Define $J(\theta)\coloneqq V^{\pi_\theta}(\rho)=\E^\pi_\rho \left[\sum^\infty_{h=0}\gamma^hr(x_h,a_h)\right]$, where $a_h\sim \pi_\theta(\cdot|x_h)$ and $x_{h+1}\sim \P(\cdot|x_h,a_h)$. \\
% To-dos
% \begin{enumerate}
%     \item Translate MDP property assumptions to $J(\theta)$ property.
%     \item Derive a performance difference $J(\theta)-J(\theta')$; regret lemma Pg36 Argarwal et. al
%     \item Projected gradient ascent
%     \item 
% \end{enumerate}
% Intended result:
% \begin{enumerate}
%     \item $L$-smooth
%     \item PL-condition on a subspace. Identify a ideal subspace $U$ s.t. $J(\theta)$ satisfies PL-condition. \\
%     % Try $D_\beta J_K(\theta)=\frac{1}{1-\gamma}\E[S^\beta_\theta(x)Q^{\pi_{\theta}}(x,\bar{a}^x_T)]$, where $S^\beta_\theta(x)=S_\theta(x)^T\beta$\\
%     Try to prove $$J^* -J(\theta)\geq\|P\|$$
% \end{enumerate}
% Assumptions:
% \begin{enumerate}
%     \item $\|\nabla_\theta f_\theta\|$
%     \item $\|\nabla_{\theta\theta} f_\theta\|$
%     \item $r(x,a)$ Lipschitz continuous in $a$, i.e. $|r(x,a)-r(x,a')|\leq L\|a-a'\|_2$ for all $a,a'\in \mathcal{A}$
% \end{enumerate}



% \section{L-smoothness of $J(\theta)$}
% % \ass For all $\theta\in \R^n$, $\|\nabla_\theta f_\theta-\nabla_\theta f_\theta'\|\leq H\|\theta-\theta'\|$ or equivalently, if $f\in C^2$, $\|\nabla_{\theta\theta}^2 f_\theta\|\leq H$\label{f lip}
% % % \lem Define $j_x(\theta)\coloneqq\E\left[r(x,\bar{a}_T^x)\right]$ 
% \lem First with result on 
% \begin{equation*}
%     \int_\mathcal{X}\rho(dx)\E\left[r(x,\bar{a}_T^x)\right]
% \end{equation*}
% consider 
% \begin{align*}
%     \nabla_\theta \int_\mathcal{X}\rho(dx)\E\left[r(x,\bar{a}_T^x)\right]
% \end{align*}
% \lem \label{j lsmooth}With \ref{f lip}, we have $J(\theta)$ is $L$-smooth, i.e. 
% \begin{equation}
%     J(\theta')\leq J(\theta)+\langle \nabla_\theta J(\theta), \theta'-\theta\rangle+\frac{L}{2}\|\theta'-\theta\|_2^2
% \end{equation}
% or equivalently if $f\in C^2$,
% \begin{equation}
%     \|\nabla_\theta J(\theta')-\nabla_\theta J(\theta)\|\leq L
% \end{equation}
% \begin{proof}

%     Let $S_\theta(x)=\int^T_0 \nabla_\theta f_\theta(t,\bar{a}^{\theta,x}_t,x)\sigma(t,\bar{a}_t^{\theta,x},x)^{-1}dB_t$  For shorthand, write $d^\theta=d^{\pi_\theta}_\rho$, $Q^{\pi_\theta}=Q^\theta$ and $\Psi(\theta,x)=\E[S_\theta(x)Q^{\pi_\theta}(x,\bar{a}_T^{\theta,x})]$
%     \begin{align}
%         \nabla_\theta J(\theta')-\nabla_\theta J(\theta)&=\frac{1}{1-\gamma}\left[\int_\mathcal{X}\E\left[S_\theta(x)Q^{\pi_\theta}(x,\bar{a}_T^{\theta,x})\right]d^{\pi_\theta}_\rho(dx)-\int_\mathcal{X}\E\left[S_{\theta'}(x)Q^{\pi_{\theta'}}(x,\bar{a}_T^{\theta,x})\right]d^{\pi_{\theta'}}_\rho(dx)\right]\\
%         &=\frac{1}{1-\gamma}\left[\int_\mathcal{X}\Psi(\theta,x)d^\theta(dx)-\int_\mathcal{X}\Psi(\theta',x)d^{\theta'}(dx)\right]\\
%         &=\frac{1}{1-\gamma}\left[\underbrace{\int_\mathcal{X}[\Psi(\theta,x)-\Psi(\theta',x)]d^{\theta}(dx)}_{I}-\underbrace{\int_\mathcal{X}\Psi(\theta',x)[d^{\theta'}-d^\theta](dx)}_{II}\right]\\
%     \end{align}

%     For $I$, we focus on the $\Psi(\theta,x)-\Psi(\theta',x)$ term. Note that under same Brownian motion coupling
%     \begin{align*}
%         \Psi(\theta,x)-\Psi(\theta',x)&=\E[S_\theta(x)Q^{\theta}(x,\bar{a}_T^{\theta,x})]-\E[S_{\theta'}(x)Q^{\theta'}(x,\bar{a}_T^{\theta,x})]\\
%         &=\E[S_\theta(x)Q^{\theta}(x,\bar{a}_T^{\theta,x})-S_{\theta'}(x)Q^{\theta'}(x,\bar{a}_T^{\theta',x})]\\
%         &=\E[S_\theta(x)Q^{\theta}(x,\bar{a}_T^{\theta,x})-S_{\theta'}(x)Q^{\theta}(x,\bar{a}_T^{\theta,x})]+\E[S_{\theta'}(x)Q^{\theta}(x,\bar{a}_T^{\theta,x})-S_{\theta'}(x)Q^{\theta'}(x,\bar{a}_T^{\theta',x})]\\
%         &\leq \|Q^{\theta}(x,\bar{a}_T^{\theta,x})\|_{\infty}E[\|S_\theta(x)-S_{\theta'}(x)\|]+\E[S_{\theta'}(x)(Q^{\theta}(x,\bar{a}_T^{\theta,x})-Q^{\theta'}(x,\bar{a}_T^{\theta',x}))]
%         % &+\E[S_{\theta'}(x)(Q^{\theta}(x,\bar{a}_T^{\theta',x})-Q^{\theta'}(x,\bar{a}_T^{\theta',x}))]
%     \end{align*}
%     \\
%     For the term 
%     \begin{align*}
%         Q^{\theta}(x,\bar{a}_T^{\theta,x})-Q^{\theta'}(x,\bar{a}_T^{\theta',x})&=r(x,\bar{a}_T^{\theta,x})-r(x,\bar{a}_T^{\theta',x})+\gamma\left(\int_\mathcal{X}V^\theta(x')P(dx'|x,\bar{a}_T^{\theta,x})-\int_\mathcal{X}V^{\theta'}(x')P(dx'|x,\bar{a}_T^{\theta',x})\right)
%     \end{align*}
%     And 
%     \begin{align*}
%         &\int_\mathcal{X}V^\theta(x')P(dx'|x,\bar{a}_T^{\theta,x})-\int_\mathcal{X}V^{\theta'}(x')P(dx'|x,\bar{a}_T^{\theta',x})\\
%         =&\int_\mathcal{X}V^\theta(x')P(dx'|x,\bar{a}_T^{\theta,x})-\int_\mathcal{X}V^{\theta'}(x')P(dx'|x,\bar{a}_T^{\theta,x})+\int_\mathcal{X}V^{\theta'}(x')P(dx'|x,\bar{a}_T^{\theta,x})-\int_\mathcal{X}V^{\theta'}(x')P(dx'|x,\bar{a}_T^{\theta',x})\\
%         =&\int_\mathcal{X}(V^\theta(x')-V^{\theta'}(x'))P(dx'|x,\bar{a}_T^{\theta,x})+\int_\mathcal{X}V^{\theta'}(x')(P(dx'|x,\bar{a}_T^{\theta,x})-P(dx'|x,\bar{a}_T^{\theta',x}))\\
%         \leq&\|V^\theta(x')-V^{\theta'}(x')\|_\infty+\int_\mathcal{X}V^{\theta'}(x')(P(dx'|x,\bar{a}_T^{\theta,x})-P(dx'|x,\bar{a}_T^{\theta',x}))\\
%         =&\|V^\theta(x')-V^{\theta'}(x')\|_\infty+\E_U\left[V^{\theta'}(\Phi_P(x,\bar{a}_T^{\theta,x},U))-V^{\theta'}(\Phi_P(x,\bar{a}_T^{\theta',x},U))\right]\\
%         \leq&\|V^\theta(x')-V^{\theta'}(x')\|_\infty+L_V\E_U\left[\|\Phi_P(x,\bar{a}_T^{\theta,x},U)-\Phi_P(x,\bar{a}_T^{\theta',x},U)\|\right]\\
%         \leq&\|V^\theta(x')-V^{\theta'}(x')\|_\infty+L_{V,x}\E_U\left[(\Phi_P(x,\bar{a}_T^{\theta,x},U)-\Phi_P(x,\bar{a}_T^{\theta',x},U))^2\right]^{1/2}\\
%         \leq&\|V^\theta(x')-V^{\theta'}(x')\|_\infty+L_{V,x}L^a_P\|\bar{a}_T^{\theta,x}-\bar{a}_T^{\theta',x}\|\\
%     \end{align*}
%     Requires stability of $V$ wrt $\theta$ and $x$ (for $\theta$, it should come from $\left\| \tilde{G}^{R}_{H}(f) - \tilde{G}^{R}_{H}(g) \right\|_{L^{2}(P)} \leq C_{3}H(C_{3}L_{K,T})^{H}d_{K,R}(f,g)$ by setting the way of parametrization); for the lipschitz wrt x, still about the same proof just setting $f=g$ but $D_0=\|x-x'\|$\\
    
%     For $II$, \\
%     We follows the similar synchronuous coupling strategy: say $(X^\theta_j,A^\theta_j)$ be the MDP process following the law introduced by $\pi_\theta$, i.e. $d\bar{a}$
    
%     , set the exogenous randomness i.i.d $\xi=(X_0,\{\zeta_j,B_j,U_j\}_{j\geq 0})$ and the recursive evolution $X^\theta_{j+1}=\Phi_P(X^\theta_j,A^\theta_j,U_j)$\\
%     The idea is to do a similar proof from Lemma B.4, with chosen way of parametrization to replace $d_{K,R}(f^\theta,f^{\theta'})$ with $\|\theta-\theta'\|$
%     \begin{align*}
%         \int_\mathcal{X}\Psi(\theta',x)[d^{\theta'}-d^\theta](dx)&=\\
%     \end{align*}
% \end{proof}


% %%--------------------------------------------------
% % \lem[Smoothness skeleton for diffusion policy gradient]
% % Let
% % \[
% % J(\theta):=V^{\pi_\theta}(\rho).
% % \]
% % Assume the policy-gradient identity holds:
% % \[
% % \nabla_\theta J(\theta)
% % =
% % \frac{1}{1-\gamma}
% % \int_{\mathcal X}
% % \Psi_\theta(x)\,d^\theta(dx),
% % \]
% % where
% % \[
% % d^\theta:=d_\rho^{\pi_\theta},
% % \]
% % and
% % \[
% % \Psi_\theta(x)
% % :=
% % \mathbb E
% % \left[
% % S_\theta^x Q_\theta(x,A_\theta^x)
% % \right].
% % \]
% % Here
% % \[
% % A_\theta^x:=\bar a_T^{\theta,x},
% % \]
% % \[
% % Q^\theta:=Q^{\pi_\theta},
% % \]
% % and
% % \[
% % S_\theta^x
% % :=
% % \int_0^T
% % \nabla_\theta f_\theta(t,\bar a_t^{\theta,x},x)
% % \sigma(t,\bar a_t^{\theta,x},x)^{-1}\,dB_t.
% % \]

% % Suppose there exist constants \(L_{\Psi,\theta},L_{\Psi,x},C_d<\infty\) such that for all
% % \(\theta,\theta'\in\Theta\) and all \(x,y\in\mathcal X\),

% % \[
% % \|\Psi_{\theta'}(x)-\Psi_\theta(x)\|
% % \le
% % L_{\Psi,\theta}\|\theta'-\theta\|,
% % \tag{A1}
% % \]

% % \[
% % \|\Psi_{\theta'}(x)-\Psi_{\theta'}(y)\|
% % \le
% % L_{\Psi,x}\|x-y\|,
% % \tag{A2}
% % \]

% % and the discounted occupancy measures satisfy the synchronous-coupling stability bound
% % \[
% % \mathbb E_{\rm sync}
% % \|X_N^{\theta'}-X_N^\theta\|
% % \le
% % C_d\|\theta'-\theta\|,
% % \tag{A3}
% % \]
% % where \(N\sim{\rm Geom}(1-\gamma)\), i.e.
% % \[
% % \mathbb P(N=j)=(1-\gamma)\gamma^j,
% % \]
% % and the coupled states satisfy
% % \[
% % X_N^\theta\sim d^\theta,\qquad X_N^{\theta'}\sim d^{\theta'}.
% % \]

% % Then \(J\) is \(L\)-smooth with
% % \[
% % L
% % =
% % \frac{L_{\Psi,\theta}+L_{\Psi,x}C_d}{1-\gamma}.
% % \]
% % That is,
% % \[
% % \|\nabla J(\theta')-\nabla J(\theta)\|
% % \le
% % L\|\theta'-\theta\|.
% % \]


% % \begin{proof}
% % By the policy-gradient formula,
% % \[
% % \nabla J(\theta')
% % -
% % \nabla J(\theta)
% % =
% % \frac{1}{1-\gamma}
% % \left[
% % \int_{\mathcal X}\Psi_{\theta'}(x)d^{\theta'}(dx)
% % -
% % \int_{\mathcal X}\Psi_\theta(x)d^\theta(dx)
% % \right].
% % \]
% % Add and subtract
% % \[
% % \int_{\mathcal X}\Psi_{\theta'}(x)d^\theta(dx)
% % \]
% % to obtain
% % \[
% % \nabla J(\theta')-\nabla J(\theta)
% % =
% % \frac{1}{1-\gamma}(I+II),
% % \]
% % where
% % \[
% % I
% % :=
% % \int_{\mathcal X}
% % \left[
% % \Psi_{\theta'}(x)-\Psi_\theta(x)
% % \right]
% % d^\theta(dx),
% % \]
% % and
% % \[
% % II
% % :=
% % \int_{\mathcal X}
% % \Psi_{\theta'}(x)
% % \left[
% % d^{\theta'}-d^\theta
% % \right](dx).
% % \]

% % For \(I\), since \(d^\theta\) is a probability measure and by Assumption \((A1)\),
% % \[
% % \begin{aligned}
% % \|I\|
% % &\le
% % \int_{\mathcal X}
% % \|\Psi_{\theta'}(x)-\Psi_\theta(x)\|
% % d^\theta(dx)
% % \\
% % &\le
% % L_{\Psi,\theta}\|\theta'-\theta\|.
% % \end{aligned}
% % \tag{1}
% % \]

% % For \(II\), use the synchronous coupling of the discounted state occupancies. Let
% % \[
% % N\sim{\rm Geom}(1-\gamma),
% % \]
% % and construct coupled trajectories so that
% % \[
% % X_N^{\theta'}\sim d^{\theta'},
% % \qquad
% % X_N^\theta\sim d^\theta.
% % \]
% % Then
% % \[
% % II
% % =
% % \mathbb E_{\rm sync}
% % \left[
% % \Psi_{\theta'}(X_N^{\theta'})
% % -
% % \Psi_{\theta'}(X_N^\theta)
% % \right].
% % \]
% % Therefore, by Assumption \((A2)\),
% % \[
% % \begin{aligned}
% % \|II\|
% % &\le
% % \mathbb E_{\rm sync}
% % \left[
% % \|\Psi_{\theta'}(X_N^{\theta'})
% % -
% % \Psi_{\theta'}(X_N^\theta)\|
% % \right]
% % \\
% % &\le
% % L_{\Psi,x}
% % \mathbb E_{\rm sync}
% % \|X_N^{\theta'}-X_N^\theta\|.
% % \end{aligned}
% % \]
% % Using Assumption \((A3)\),
% % \[
% % \|II\|
% % \le
% % L_{\Psi,x}C_d\|\theta'-\theta\|.
% % \tag{2}
% % \]

% % Combining \((1)\) and \((2)\), we get
% % \[
% % \begin{aligned}
% % \|\nabla J(\theta')-\nabla J(\theta)\|
% % &\le
% % \frac{1}{1-\gamma}
% % \left(
% % \|I\|+\|II\|
% % \right)
% % \\
% % &\le
% % \frac{L_{\Psi,\theta}+L_{\Psi,x}C_d}{1-\gamma}
% % \|\theta'-\theta\|.
% % \end{aligned}
% % \]
% % Thus \(J\) is \(L\)-smooth with
% % \[
% % L=
% % \frac{L_{\Psi,\theta}+L_{\Psi,x}C_d}{1-\gamma}.
% % \]
% % \end{proof}

% % \lem[\(\Psi\)-Lipschitz]
% % Assume that for all \(x,x'\in\mathcal X\) and \(\theta,\theta'\in\Theta\), under the same-Brownian coupling,
% % \[
% % \left(
% % \mathbb E
% % \|A_{\theta'}^{x'}-A_\theta^x\|^2
% % \right)^{1/2}
% % \le
% % C_{A,x}\|x'-x\|+C_{A,\theta}\|\theta'-\theta\|,
% % \tag{B1}
% % \]
% % and
% % \[
% % \left(
% % \mathbb E
% % \|S_{\theta'}^{x'}-S_\theta^x\|^2
% % \right)^{1/2}
% % \le
% % C_{S,x}\|x'-x\|+C_{S,\theta}\|\theta'-\theta\|.
% % \tag{B2}
% % \]
% % Assume also
% % \[
% % \left(
% % \mathbb E\|S_\theta^x\|^2
% % \right)^{1/2}
% % \le M_S,
% % \qquad
% % \|Q_\theta\|_\infty\le M_Q,
% % \tag{B3}
% % \]
% % and
% % \[
% % |Q_{\theta'}(x',a')-Q_\theta(x,a)|
% % \le
% % L_Q^x\|x'-x\|
% % +
% % L_Q^a\|a'-a\|
% % +
% % C_Q^\theta\|\theta'-\theta\|.
% % \tag{B4}
% % \]
% % Then
% % \[
% % \|\Psi_{\theta'}(x')-\Psi_\theta(x)\|
% % \le
% % L_{\Psi,x}\|x'-x\|
% % +
% % L_{\Psi,\theta}\|\theta'-\theta\|,
% % \]
% % where one may take
% % \[
% % L_{\Psi,x}
% % =
% % M_QC_{S,x}
% % +
% % M_S(L_Q^x+L_Q^aC_{A,x}),
% % \]
% % and
% % \[
% % L_{\Psi,\theta}
% % =
% % M_QC_{S,\theta}
% % +
% % M_S(L_Q^aC_{A,\theta}+C_Q^\theta).
% % \]


% % \begin{proof}
% % Under the same-Brownian coupling,
% % \[
% % \begin{aligned}
% % &\Psi_{\theta'}(x')-\Psi_\theta(x)\\
% % % &=\Psi_{\theta'}(x')-\Psi_{\theta'}(x)+\Psi_{\theta'}(x)-\Psi_{\theta}(x')\\
% % &=
% % \mathbb E
% % \left[
% % S_{\theta'}^{x'}Q_{\theta'}(x',A_{\theta'}^{x'})
% % -
% % S_\theta^xQ_\theta(x,A_\theta^x)
% % \right]
% % \\
% % &=
% % \mathbb E
% % \left[
% % (S_{\theta'}^{x'}-S_\theta^x)Q_\theta(x,A_\theta^x)
% % \right]
% % \\
% % &\quad+
% % \mathbb E
% % \left[
% % S_{\theta'}^{x'}
% % \left(
% % Q_{\theta'}(x',A_{\theta'}^{x'})
% % -
% % Q_\theta(x,A_\theta^x)
% % \right)
% % \right].
% % \end{aligned}
% % \]
% % Hence, by Cauchy--Schwarz,
% % \[
% % \begin{aligned}
% % \|\Psi_{\theta'}(x')-\Psi_\theta(x)\|
% % &\le
% % M_Q
% % \left(
% % \mathbb E
% % \|S_{\theta'}^{x'}-S_\theta^x\|^2
% % \right)^{1/2}
% % \\
% % &\quad+
% % M_S
% % \left(
% % \mathbb E
% % \left|
% % Q_{\theta'}(x',A_{\theta'}^{x'})
% % -
% % Q_\theta(x,A_\theta^x)
% % \right|^2
% % \right)^{1/2}.
% % \end{aligned}
% % \]
% % Using \((B1)\), \((B2)\), and \((B4)\) gives
% % \[
% % \|\Psi_{\theta'}(x')-\Psi_\theta(x)\|
% % \le
% % \left[
% % M_QC_{S,x}
% % +
% % M_S(L_Q^x+L_Q^aC_{A,x})
% % \right]\|x'-x\|
% % \]
% % \[
% % +
% % \left[
% % M_QC_{S,\theta}
% % +
% % M_S(L_Q^aC_{A,\theta}+C_Q^\theta)
% % \right]\|\theta'-\theta\|.
% % \]
% % This proves the claim.
% % \end{proof}

% % \lem[Discounted occupancy stability via synchronous coupling]
% % Assume that the coupled action processes satisfy
% % \[
% % E_j
% % :=
% % \left(
% % \mathbb E_{\rm sync}
% % \|A_j^{\theta'}-A_j^\theta\|^2
% % \right)^{1/2}
% % \le
% % C_A^xD_j+C_A^\theta\|\theta'-\theta\|,
% % \tag{C1}
% % \]
% % where
% % \[
% % D_j
% % :=
% % \left(
% % \mathbb E_{\rm sync}
% % \|X_j^{\theta'}-X_j^\theta\|^2
% % \right)^{1/2}.
% % \]
% % Assume the transition kernel admits a synchronous \(L^2\)-coupling:
% % \[
% % \left(
% % \mathbb E_U
% % \|\Phi_P(x,a,U)-\Phi_P(x',a',U)\|^2
% % \right)^{1/2}
% % \le
% % L_P^x\|x-x'\|+L_P^a\|a-a'\|.
% % \tag{C2}
% % \]
% % Let
% % \[
% % \alpha:=L_P^x+L_P^aC_A^x,
% % \qquad
% % \beta:=L_P^aC_A^\theta.
% % \]
% % If
% % \[
% % \gamma\alpha<1,
% % \]
% % then
% % \[
% % \mathbb E_{\rm sync}
% % \|X_N^{\theta'}-X_N^\theta\|
% % \le
% % \frac{\gamma\beta}{1-\gamma\alpha}
% % \|\theta'-\theta\|.
% % \]
% % Thus one may take
% % \[
% % C_d
% % =
% % \frac{\gamma L_P^aC_A^\theta}
% % {1-\gamma(L_P^x+L_P^aC_A^x)}.
% % \]

% % \begin{proof}
% % Construct the coupled trajectories using the same initial state, same Brownian motions, same initial diffusion noises, and same transition seeds:
% % \[
% % X_0^{\theta'}=X_0^\theta,
% % \]
% % \[
% % A_j^\theta
% % =
% % \bar a_T^{\theta,X_j^\theta},
% % \qquad
% % A_j^{\theta'}
% % =
% % \bar a_T^{\theta',X_j^{\theta'}},
% % \]
% % and
% % \[
% % X_{j+1}^\theta
% % =
% % \Phi_P(X_j^\theta,A_j^\theta,U_j),
% % \qquad
% % X_{j+1}^{\theta'}
% % =
% % \Phi_P(X_j^{\theta'},A_j^{\theta'},U_j).
% % \]
% % By \((C2)\),
% % \[
% % D_{j+1}
% % \le
% % L_P^xD_j+L_P^aE_j.
% % \]
% % Using \((C1)\),
% % \[
% % D_{j+1}
% % \le
% % (L_P^x+L_P^aC_A^x)D_j
% % +
% % L_P^aC_A^\theta\|\theta'-\theta\|.
% % \]
% % Thus
% % \[
% % D_{j+1}
% % \le
% % \alpha D_j+\beta\|\theta'-\theta\|.
% % \]
% % Since \(D_0=0\),
% % \[
% % D_j
% % \le
% % \beta
% % \sum_{k=0}^{j-1}\alpha^k
% % \|\theta'-\theta\|.
% % \]
% % Now let \(N\sim{\rm Geom}(1-\gamma)\). Then
% % \[
% % \begin{aligned}
% % \mathbb E_{\rm sync}
% % \|X_N^{\theta'}-X_N^\theta\|
% % &=
% % (1-\gamma)
% % \sum_{j=0}^\infty
% % \gamma^j
% % \mathbb E_{\rm sync}
% % \|X_j^{\theta'}-X_j^\theta\|
% % \\
% % &\le
% % (1-\gamma)
% % \sum_{j=0}^\infty
% % \gamma^jD_j
% % \\
% % &\le
% % (1-\gamma)
% % \sum_{j=0}^\infty
% % \gamma^j
% % \beta
% % \sum_{k=0}^{j-1}\alpha^k
% % \|\theta'-\theta\|.
% % \end{aligned}
% % \]
% % Since
% % \[
% % (1-\gamma)
% % \sum_{j=0}^\infty
% % \gamma^j
% % \sum_{k=0}^{j-1}\alpha^k
% % =
% % \frac{\gamma}{1-\gamma\alpha},
% % \]
% % provided \(\gamma\alpha<1\), we obtain
% % \[
% % \mathbb E_{\rm sync}
% % \|X_N^{\theta'}-X_N^\theta\|
% % \le
% % \frac{\gamma\beta}{1-\gamma\alpha}
% % \|\theta'-\theta\|.
% % \]
% % \end{proof}
% %-----------------------------
\newpage
\section{Smoothness of the diffusion-policy objective}

Let
\[
J(\theta):=V^{\pi_\theta}(\rho).
\]
For the discounted occupancy measure, use the paper's notation
\[
d_\rho^{\pi_\theta}(dx)
:=
\int_{\mathcal X}d^{\pi_\theta}(dx\mid x')\,\rho(dx'),
\]
where
\[
d^{\pi_\theta}(dx\mid x')
:=
(1-\gamma)\sum_{n=0}^{\infty}\gamma^n
P_{\pi_\theta}^n(dx\mid x').
\]
Here
\[
P_{\pi_\theta}(dy\mid x)
:=
\int_{\mathcal A}P(dy\mid x,a)\pi_\theta(da\mid x)
\]
is the policy-induced state transition kernel.

For fixed \(x\), let
\[
d\bar a_t^{\theta,x}
=
f_\theta(t,\bar a_t^{\theta,x},x)\,dt
+
\sigma(t,\bar a_t^{\theta,x},x)\,dB_t,
\qquad
\bar a_0^{\theta,x}=\zeta\sim\mu_0.
\]
Set
\[
A_\theta^x:=\bar a_T^{\theta,x}.
\]
The Girsanov score is
\[
S_\theta^x
:=
\int_0^T
\nabla_\theta f_\theta(t,\bar a_t^{\theta,x},x)
\sigma(t,\bar a_t^{\theta,x},x)^{-1}\,dB_t.
\]
Define
\[
Q_\theta:=Q^{\pi_\theta},
\qquad
V_\theta:=V^{\pi_\theta},
\]
and
\[
\Psi_\theta(x)
:=
\mathbb E
\left[
S_\theta^x Q_\theta(x,A_\theta^x)
\right].
\]
By Proposition 5.2,
\[
\nabla J(\theta)
=
\frac{1}{1-\gamma}
\int_{\mathcal X}
\Psi_\theta(x)\,d_\rho^{\pi_\theta}(dx).
\tag{PG}
\]

\subsection{Additional smoothness assumptions}

\begin{ass}[Strong common-noise realization]
For every \((\theta,x)\), the action SDE admits a unique strong solution. Hence, on one common stochastic basis carrying \((\zeta,B)\), we may realize
\[
\bar a_\cdot^{\theta,x}
=
\mathcal I_{\theta,x}(\zeta,B)
\]
for all relevant \((\theta,x)\).
\end{ass}

\begin{ass}[Ellipticity requirement]
For all $t,x,a$, $\sigma(t,x,a)$ is symmetric and 
\begin{equation}
     \kappa I \preceq \sigma(t,a,x)\sigma(t,a,x)^T\preceq \Lambda I
\end{equation}
\end{ass}

\begin{ass}[moment bound of initial distribution]
Assume $\E\|\zeta\|^4<\infty$
\end{ass}

\begin{ass}[Action-SDE Lipschitz regularity]
% There exist constants \(L_f^a,L_f^x,L_f^\theta,L_\sigma^a,L_\sigma^x<\infty\) such that
% \[
% \|f_{\theta'}(t,a',x')-f_\theta(t,a,x)\|
% \le
% L_f^a\|a'-a\|
% +
% L_f^x\|x'-x\|
% +
% L_f^\theta\|\theta'-\theta\|,
% \]
% and
% \[
% \|\sigma(t,a',x')-\sigma(t,a,x)\|_{\rm HS}
% \le
% L_\sigma^a\|a'-a\|
% +
% L_\sigma^x\|x'-x\|.
% \]
There exist constants \(L_f^a,L_f^x,L_\sigma^a,L_\sigma^x<\infty\) such that
\[
\|f_{\theta}(t,a',x')-f_\theta(t,a,x)\|
\le
L_f^a\|a'-a\|
+
L_f^x\|x'-x\|
,
\]
and
\[
\|\sigma(t,a',x')-\sigma(t,a,x)\|_{\rm HS}
\le
L_\sigma^a\|a'-a\|
+
L_\sigma^x\|x'-x\|.
\]
\end{ass}


% \ass[Score-integrand regularity]
% Define
% \[
% G_\theta(t,a,x)
% :=
% \nabla_\theta f_\theta(t,a,x)\sigma(t,a,x)^{-1}.
% \]
% Assume
% \[
% \|G_\theta(t,a,x)\|_{\rm HS}\le M_G,
% \]
% and
% \[
% \|G_{\theta'}(t,a',x')-G_\theta(t,a,x)\|_{\rm HS}
% \le
% L_G^\theta\|\theta'-\theta\|
% +
% L_G^a\|a'-a\|
% +
% L_G^x\|x'-x\|.
% \]
\begin{ass}[$f$ bound]
Assume uniformly in \((t,\theta,x)\),
\[
\|f_\theta(t,a,x)\|
\le B_f(1+\|a\|),
\]
\end{ass}
\begin{ass}[$\nabla f$ regularity]
Assume $\nabla f$ satisfies
\[
\|\nabla_\theta f_\theta(t,a,x)\|_{HS}\leq B_{\nabla f}*(1+\|a\|)
\]
and $\Theta$ convex
\end{ass}
% and 
% \[C^{-1}I_d\prec\sigma(t,a,x)\prec CI_d\]
% that is $\|\sigma(s,a,x)\|_{\rm op}\leq C$
\begin{ass}[$\nabla_\theta f$ Lipschitz]
There exist constants \(L_{\nabla f}^a,L_{\nabla f}^x,L_{\nabla f}^\theta\) $<\infty$ such that
\[
\|\nabla_\theta f_{\theta'}(t,a',x')-\nabla_\theta f_\theta(t,a,x)\|_{\rm HS}
\le
L_{\nabla f}^a\|a'-a\|
+
L_{\nabla f}^x\|x'-x\|
+
L_{\nabla f}^\theta\|\theta'-\theta\|,
\]
\end{ass}

% \begin{ass}[\(\mathcal F_{K,s}\)-compatible parametrization] (Replacable by the uniform bound of $\nabla f$?)

% \textcolor{purple}{Let \(\Theta\subset\mathbb R^p\). For every \(\theta\in\Theta\),
% \[
% f_\theta\in\mathcal F_{K,s}.
% \]
% Moreover, for every localization radius \(R\ge1\), there exists a constant
% \(L_{\Theta,R}<\infty\) such that
% \[
% d_{K,R}(f_\theta,f_{\theta'})
% \le
% L_{\Theta,R}\|\theta-\theta'\|,
% \qquad
% \forall \theta,\theta'\in\Theta,
% \]
% where
% \[
% d_{K,R}(f,g)
% :=
% \sup_{(t,a,x)\in[0,T]\times B_R^d\times[0,1]^m}
% \frac{\|f(t,a,x)-g(t,a,x)\|}{K}.
% \]}
% \end{ass}


\begin{ass}[Environment regularity]
Assume the reward satisfies the paper's Assumption 4.2:
\[
|r(x,a)-r(x',a')|
\le
L_r^x\|x-x'\|+L_r^a\|a-a'\|.
\]
Assume also \(|r(x,a)|\le r_{\max}\). The transition kernel admits the synchronous generative representation
\[
\Phi_P(x,a,U)\sim P(\cdot\mid x,a),
\qquad
U\sim{\rm Unif}[0,1],
\]
and
\[
\left(
\mathbb E_U
\|\Phi_P(x,a,U)-\Phi_P(x',a',U)\|^2
\right)^{1/2}
\le
L_P^x\|x-x'\|+L_P^a\|a-a'\|.
\tag{E}
\]
\end{ass}

\subsection{Property lemmas}
% \begin{lem}[Finite moments of $\bar{a}_t^{\theta,x}$]
%     Assume $\bar{a}_0^{\theta,x}=\zeta$ with $\E\|\zeta\|^2 <\infty.$ Then we have $\sup_{t,\theta,x}\E\|\bar{a}_t^{\theta,x}\|<\infty$

%     \begin{proof}
%         By Ito's formula, $\|\bar{a}_t^{\theta,x}\|$
%     \end{proof}
% \end{lem}
\begin{lem}[Uniform second and fourth moments of
\(\bar a_t^{\theta,x}\)]
Assume that
\[
\bar a_0^{\theta,x}=\zeta,
\]

If \(\mathbb E\|\zeta\|^2<\infty\), then
\[
M_{A,2}
:=
\sup_{\theta,x}\sup_{0\le t\le T}
\mathbb E\|\bar a_t^{\theta,x}\|^2
<\infty,
\]
and
\[
M_{A,2}
\le
e^{c_2T}\mathbb E\|\zeta\|^2
+
\frac{b_2}{c_2}\left(e^{c_2T}-1\right),
\]
where
\[
c_2:=1+2B_f^2,
\qquad
b_2:=2B_f^2+d_a\Lambda.
\]

If, in addition, \(\mathbb E\|\zeta\|^4<\infty\), then
\[
M_{A,4}
:=
\sup_{\theta,x}\sup_{0\le t\le T}
\mathbb E\|\bar a_t^{\theta,x}\|^4
<\infty,
\]
and
\[
M_{A,4}
\le
e^{c_4T}\mathbb E\|\zeta\|^4
+
\frac{b_4}{c_4}\left(e^{c_4T}-1\right),
\]
where
\[
c_4:=8B_f+2(d_a+2)\Lambda,
\qquad
b_4:=4B_f+2(d_a+2)\Lambda.
\]

\begin{proof}
Fix \((\theta,x)\) and write
\[
\bar a_t:=\bar a_t^{\theta,x},
\qquad
f_t:=f_\theta(t,\bar a_t,x),
\qquad
\sigma_t:=\sigma(t,\bar a_t,x).
\]

We first establish the second-moment estimate. It\^o's
formula gives
\[
\begin{aligned}
d\|\bar a_t\|^2
={}&
\left[
2\langle\bar a_t,f_t\rangle
+
\|\sigma_t\|_{\rm HS}^2
\right]dt
+
2\langle\bar a_t,\sigma_t\,dB_t\rangle.
\end{aligned}
\]
By Young's inequality and the linear-growth assumption,
\[
\begin{aligned}
2\langle a,f_\theta(t,a,x)\rangle
&\le
\|a\|^2+\|f_\theta(t,a,x)\|^2\\
&\le
\|a\|^2+B_f^2(1+\|a\|)^2\\
&\le
(1+2B_f^2)\|a\|^2+2B_f^2.
\end{aligned}
\]
Moreover,
\[
\|\sigma(t,a,x)\|_{\rm HS}^2
=
\operatorname{tr}\!\left(
\sigma(t,a,x)\sigma(t,a,x)^\top
\right)
\le d_a\Lambda.
\]
After taking expectations (For full rigorous proof, use stopping process $\bar{a}_{t\wedge\tau_n} $ with $\tau_n=\min\{t:\|\bar{a}_t\|\geq n\}$. Pushing $n\to \infty$, Fatou's lemma gives $\int^t_0 \|\sigma_s^T\bar{a}_s\|^2ds<\infty$ and $\int^t_0 \langle\bar{a}_s,\sigma_sdB_s\rangle$ is a martingale), we therefore obtain
\[
\frac{d}{dt}\mathbb E\|\bar a_t\|^2
\le
c_2\mathbb E\|\bar a_t\|^2+b_2.
\]
Gr\"onwall's inequality yields
\[
\mathbb E\|\bar a_t\|^2
\le
e^{c_2t}\mathbb E\|\zeta\|^2
+
\frac{b_2}{c_2}\left(e^{c_2t}-1\right).
\]
Taking the supremum over \((\theta,x)\) and \(t\in[0,T]\) proves the
second-moment estimate.

For the fourth moment, apply It\^o's formula to
\[
\phi(a):=\|a\|^4.
›\]
Since
\[
\nabla\phi(a)=4\|a\|^2a,
\qquad
\nabla^2\phi(a)
=
8aa^\top+4\|a\|^2I_{d_a},
\]
we have
\[
\begin{aligned}
d\|\bar a_t\|^4
={}&
\Big[
4\|\bar a_t\|^2\langle\bar a_t,f_t\rangle
+
4\|\sigma_t^\top\bar a_t\|^2\\
&\hspace{23mm}
+
2\|\bar a_t\|^2\|\sigma_t\|_{\rm HS}^2
\Big]dt\\
&+
4\|\bar a_t\|^2
\langle\bar a_t,\sigma_t\,dB_t\rangle.
\end{aligned}
\tag{1}
\]

For the drift contribution,
\[
\begin{aligned}
4\|a\|^2\langle a,f_\theta(t,a,x)\rangle
&\le
4\|a\|^3\|f_\theta(t,a,x)\|\\
&\le
4B_f\|a\|^3(1+\|a\|)\\
&\le
4B_f+8B_f\|a\|^4,
\end{aligned}
\tag{2}
\]
where we used \(\|a\|^3\le 1+\|a\|^4\).

The upper ellipticity bound implies
\[
\|\sigma_t^\top a\|^2
\le \Lambda\|a\|^2,
\qquad
\|\sigma_t\|_{\rm HS}^2
\le d_a\Lambda.
\]
Consequently,
\[
\begin{aligned}
4\|\sigma_t^\top a\|^2
+
2\|a\|^2\|\sigma_t\|_{\rm HS}^2
&\le
2(d_a+2)\Lambda\|a\|^2\\
&\le
2(d_a+2)\Lambda(1+\|a\|^4).
\end{aligned}
\tag{3}
\]
Combining \((1)\)--\((3)\) and taking expectations gives
\[
\frac{d}{dt}\mathbb E\|\bar a_t\|^4
\le
c_4\mathbb E\|\bar a_t\|^4+b_4.
\]
Another application of Gr\"onwall's inequality yields
\[
\mathbb E\|\bar a_t\|^4
\le
e^{c_4t}\mathbb E\|\zeta\|^4
+
\frac{b_4}{c_4}\left(e^{c_4t}-1\right).
\]
Taking the supremum over \((\theta,x)\) and \(t\in[0,T]\) proves the
fourth-moment estimate.

% The preceding calculations can be justified rigorously by first
% applying It\^o's formula to the stopped process
% \[
% \bar a_{t\wedge\tau_n},
% \qquad
% \tau_n:=\inf\{t\ge0:\|\bar a_t\|\ge n\},
% \]
% and then letting \(n\to\infty\).
\end{proof}
\end{lem}

\begin{lem}[Alternative of assumption 1 ]
Assume $\zeta\sim \mathcal{N}(0,\sigma_\zeta^2)$
    
\end{lem}
\begin{lem}[Lipschitz of $f$ w.r.t. $\theta$]
Assume $\Theta$ is convex. For $\theta,\theta'\in \Theta$, we have
\begin{equation}
    \|f_{\theta'}(t,a,x)-f_\theta(t,a,x)\|\leq B_{\nabla f}(1+\|a\|)\|\theta'-\theta\|
\end{equation}
\begin{proof}
    By the regularity assumption of $\nabla f$, 
    \[\|f(t,a,x,\theta')-f(t,a,x,\theta)\|=\left\|\int^1_0\nabla_\theta f (t,a,x, (1-s)\theta+s\theta')(\theta'-\theta)ds\right\| \leq B_{\nabla f}(1+\|a\|)\|\theta'-\theta\|\]
\end{proof}
\end{lem}

\begin{lem}[Action-path stability]
Under the same-initial-noise and same-Brownian coupling, there exist
constants \(C_A^x,C_A^\theta<\infty\) such that
\[
\sup_{0\le t\le T}
\left(
\mathbb E
\|\bar a_t^{\theta',x'}-\bar a_t^{\theta,x}\|^2
\right)^{1/2}
\le
C_A^x\|x'-x\|
+
C_A^\theta\|\theta'-\theta\|.
\tag{A}
\]
In particular, since
\[
A_\theta^x=\bar a_T^{\theta,x},
\]
we have
\[
\left(
\mathbb E
\|A_{\theta'}^{x'}-A_\theta^x\|^2
\right)^{1/2}
\le
C_A^x\|x'-x\|
+
C_A^\theta\|\theta'-\theta\|.
\tag{A'}
\]

\begin{proof}
Let
\[
Z_t
:=
\bar a_t^{\theta',x'}-\bar a_t^{\theta,x}.
\]
Because the two processes use the same initial random variable,
\[
Z_0=0.
\]

Under the same-Brownian coupling,
\[
\begin{aligned}
dZ_t
={}&
\left[
f_{\theta'}(t,\bar a_t^{\theta',x'},x')
-
f_\theta(t,\bar a_t^{\theta,x},x)
\right]dt\\
&+
\left[
\sigma(t,\bar a_t^{\theta',x'},x')
-
\sigma(t,\bar a_t^{\theta,x},x)
\right]dB_t.
\end{aligned}
\tag{1}
\]

By the Lipschitz assumptions on \(f_\theta\), convexity of
\(\Theta\), and the fundamental theorem of calculus,
\[
\begin{aligned}
&
\left\|
f_{\theta'}(t,\bar a_t^{\theta',x'},x')
-
f_\theta(t,\bar a_t^{\theta,x},x)
\right\|\\
&\le
L_f^a\|Z_t\|
+
L_f^x\|x'-x\|\\
&\quad+
B_{\nabla f}
\left(1+\|\bar a_t^{\theta,x}\|\right)
\|\theta'-\theta\|.
\end{aligned}
\tag{2}
\]
Similarly,
\[
\begin{aligned}
&
\left\|
\sigma(t,\bar a_t^{\theta',x'},x')
-
\sigma(t,\bar a_t^{\theta,x},x)
\right\|_{\rm HS}\\
&\le
L_\sigma^a\|Z_t\|
+
L_\sigma^x\|x'-x\|.
\end{aligned}
\tag{3}
\]

Applying It\^o's formula to \(\|Z_t\|^2\) gives
\[
\begin{aligned}
d\|Z_t\|^2
={}&
2\left\langle
Z_t,
f_{\theta'}(t,\bar a_t^{\theta',x'},x')
-
f_\theta(t,\bar a_t^{\theta,x},x)
\right\rangle dt\\
&+
\left\|
\sigma(t,\bar a_t^{\theta',x'},x')
-
\sigma(t,\bar a_t^{\theta,x},x)
\right\|_{\rm HS}^2dt\\
&+
2\left\langle
Z_t,
\left[
\sigma(t,\bar a_t^{\theta',x'},x')
-
\sigma(t,\bar a_t^{\theta,x},x)
\right]dB_t
\right\rangle.
\end{aligned}
\tag{4}
\]

Using \((2)\) and Young's inequality,
\[
\begin{aligned}
&
2\left\langle
Z_t,
f_{\theta'}(t,\bar a_t^{\theta',x'},x')
-
f_\theta(t,\bar a_t^{\theta,x},x)
\right\rangle\\
&\le
2L_f^a\|Z_t\|^2
+
2L_f^x\|Z_t\|\|x'-x\|\\
&\quad+
2B_{\nabla f}
\left(1+\|\bar a_t^{\theta,x}\|\right)
\|Z_t\|\|\theta'-\theta\|\\
&\le
(2L_f^a+1)\|Z_t\|^2
+
2(L_f^x)^2\|x'-x\|^2\\
&\quad+
2B_{\nabla f}^2
\left(1+\|\bar a_t^{\theta,x}\|\right)^2
\|\theta'-\theta\|^2.
\end{aligned}
\tag{5}
\]
Moreover, \((3)\) implies
\[
\begin{aligned}
&
\left\|
\sigma(t,\bar a_t^{\theta',x'},x')
-
\sigma(t,\bar a_t^{\theta,x},x)
\right\|_{\rm HS}^2\\
&\le
2(L_\sigma^a)^2\|Z_t\|^2
+
2(L_\sigma^x)^2\|x'-x\|^2.
\end{aligned}
\tag{6}
\]

Taking expectations in \((4)\), using that the stochastic integral has
mean zero, and observing that
\[
\begin{aligned}
\mathbb E
\left(1+\|\bar a_t^{\theta,x}\|\right)^2
&\le
2\left(
1+\mathbb E\|\bar a_t^{\theta,x}\|^2
\right)\\
&\le
2(1+M_{A,2}),
\end{aligned}
\]
we obtain
\[
\begin{aligned}
\frac{d}{dt}\mathbb E\|Z_t\|^2
\le{}&
c_{Z,2}\mathbb E\|Z_t\|^2
+
b_{Z,2}^x\|x'-x\|^2\\
&+
b_{Z,2}^\theta\|\theta'-\theta\|^2,
\end{aligned}
\tag{7}
\]
where
\[
\begin{aligned}
c_{Z,2}
&:=
2L_f^a+1+2(L_\sigma^a)^2,\\
b_{Z,2}^x
&:=
2(L_f^x)^2+2(L_\sigma^x)^2,\\
b_{Z,2}^\theta
&:=
4B_{\nabla f}^2(1+M_{A,2}).
\end{aligned}
\tag{8}
\]

Since \(Z_0=0\), Gr\"onwall's inequality yields
\[
\begin{aligned}
\mathbb E\|Z_t\|^2
\le
\frac{e^{c_{Z,2}t}-1}{c_{Z,2}}
\left[
b_{Z,2}^x\|x'-x\|^2
+
b_{Z,2}^\theta\|\theta'-\theta\|^2
\right].
\end{aligned}
\]
Consequently,
\[
\begin{aligned}
\left(\mathbb E\|Z_t\|^2\right)^{1/2}
\le{}&
\left[
\frac{e^{c_{Z,2}t}-1}{c_{Z,2}}
b_{Z,2}^x
\right]^{1/2}
\|x'-x\|\\
&+
\left[
\frac{e^{c_{Z,2}t}-1}{c_{Z,2}}
b_{Z,2}^\theta
\right]^{1/2}
\|\theta'-\theta\|.
\end{aligned}
\]

Therefore, setting
\[
C_A^x
:=
\left[
\frac{e^{c_{Z,2}T}-1}{c_{Z,2}}
b_{Z,2}^x
\right]^{1/2},
\qquad
C_A^\theta
:=
\left[
\frac{e^{c_{Z,2}T}-1}{c_{Z,2}}
b_{Z,2}^\theta
\right]^{1/2},
\]
and taking the supremum over \(t\in[0,T]\) proves \((A)\).
Evaluating the estimate at \(t=T\) proves \((A')\).
\end{proof}
\end{lem}

\begin{lem}[Fourth-moment action-path stability]
Under the same-initial-noise and same-Brownian coupling, there exist
constants \(C_{A,4}^x,C_{A,4}^\theta<\infty\) such that
\[
\sup_{0\le t\le T}
\left(
\mathbb E
\|\bar a_t^{\theta',x'}-\bar a_t^{\theta,x}\|^4
\right)^{1/4}
\le
C_{A,4}^x\|x'-x\|
+
C_{A,4}^\theta\|\theta'-\theta\|.
\tag{A4}
\]
In particular,
\[
\left(
\mathbb E
\|A_{\theta'}^{x'}-A_\theta^x\|^4
\right)^{1/4}
\le
C_{A,4}^x\|x'-x\|
+
C_{A,4}^\theta\|\theta'-\theta\|.
\tag{A4'}
\]

\begin{proof}
Define $Z_t$ similarly to Lemma 1.3
% Let
% \[
% Z_t
% :=
% \bar a_t^{\theta',x'}-\bar a_t^{\theta,x}.
% \]
% Because the two processes use the same initial random variable,
% \[
% Z_0=0.
% \]
% Under the same-Brownian coupling,
% \[
% \begin{aligned}
% dZ_t
% ={}&
% \left[
% f_{\theta'}(t,\bar a_t^{\theta',x'},x')
% -
% f_\theta(t,\bar a_t^{\theta,x},x)
% \right]dt\\
% &+
% \left[
% \sigma(t,\bar a_t^{\theta',x'},x')
% -
% \sigma(t,\bar a_t^{\theta,x},x)
% \right]dB_t.
% \end{aligned}
% \tag{1}
% \]

% The drift assumptions, convexity of the parameter domain, and the
% fundamental theorem of calculus imply
% \[
% \begin{aligned}
% &
% \left\|
% f_{\theta'}(t,\bar a_t^{\theta',x'},x')
% -
% f_\theta(t,\bar a_t^{\theta,x},x)
% \right\|\\
% &\le
% L_f^a\|Z_t\|
% +
% L_f^x\|x'-x\|\\
% &\quad+
% B_{\nabla f}
% \left(1+\|\bar a_t^{\theta,x}\|\right)
% \|\theta'-\theta\|.
% \end{aligned}
% \tag{2}
% \]
% Similarly,
% \[
% \begin{aligned}
% &
% \left\|
% \sigma(t,\bar a_t^{\theta',x'},x')
% -
% \sigma(t,\bar a_t^{\theta,x},x)
% \right\|_{\rm HS}\\
% &\le
% L_\sigma^a\|Z_t\|
% +
% L_\sigma^x\|x'-x\|.
% \end{aligned}
% \tag{3}
% \]

For an SDE \(dZ_t=F_tdt+\Sigma_t dB_t\), It\^o's formula gives
\[
\begin{aligned}
d\|Z_t\|^4
={}&
\left[
4\|Z_t\|^2\langle Z_t,F_t\rangle
+
2\|Z_t\|^2\|\Sigma_t\|_{\rm HS}^2
+
4\|\Sigma_t^\top Z_t\|^2
\right]dt\\
&+
4\|Z_t\|^2
\langle Z_t,\Sigma_t\,dB_t\rangle.
\end{aligned}
\tag{4}
\]
Since
\[
\|\Sigma_t^\top Z_t\|^2
\le
\|Z_t\|^2\|\Sigma_t\|_{\rm HS}^2,
\]
the two diffusion correction terms satisfy
\[
2\|Z_t\|^2\|\Sigma_t\|_{\rm HS}^2
+
4\|\Sigma_t^\top Z_t\|^2
\le
6\|Z_t\|^2\|\Sigma_t\|_{\rm HS}^2.
\tag{5}
\]

By \((2)\),
\[
\begin{aligned}
4\|Z_t\|^2\langle Z_t,F_t\rangle
\le{}&
4L_f^a\|Z_t\|^4
+
4L_f^x\|Z_t\|^3\|x'-x\|\\
&+
4B_{\nabla f}
\left(1+\|\bar a_t^{\theta,x}\|\right)
\|Z_t\|^3\|\theta'-\theta\|.
\end{aligned}
\]
Using
\[
4u^3v\le 3u^4+v^4,
\qquad u,v\ge0,
\]
we obtain
\[
\begin{aligned}
4\|Z_t\|^2\langle Z_t,F_t\rangle
\le{}&
(4L_f^a+6)\|Z_t\|^4
+
(L_f^x)^4\|x'-x\|^4\\
&+
B_{\nabla f}^4
\left(1+\|\bar a_t^{\theta,x}\|\right)^4
\|\theta'-\theta\|^4.
\end{aligned}
\tag{6}
\]

On the other hand, \((3)\) gives
\[
\begin{aligned}
6\|Z_t\|^2\|\Sigma_t\|_{\rm HS}^2
\le{}&
12(L_\sigma^a)^2\|Z_t\|^4\\
&+
12(L_\sigma^x)^2
\|Z_t\|^2\|x'-x\|^2.
\end{aligned}
\]
Since
\[
2\|Z_t\|^2\|x'-x\|^2
\le
\|Z_t\|^4+\|x'-x\|^4,
\]
it follows that
\[
\begin{aligned}
6\|Z_t\|^2\|\Sigma_t\|_{\rm HS}^2
\le{}&
\left[
12(L_\sigma^a)^2
+
6(L_\sigma^x)^2
\right]\|Z_t\|^4\\
&+
6(L_\sigma^x)^2\|x'-x\|^4.
\end{aligned}
\tag{7}
\]

Taking expectations in \((4)\), using that the stochastic integral has
mean zero, and combining \((5)\)--\((7)\), we obtain
\[
\begin{aligned}
\frac{d}{dt}\mathbb E\|Z_t\|^4
\le{}&
c_{Z,4}\mathbb E\|Z_t\|^4
+
b_{Z,4}^x\|x'-x\|^4\\
&+
B_{\nabla f}^4
\mathbb E
\left(1+\|\bar a_t^{\theta,x}\|\right)^4
\|\theta'-\theta\|^4,
\end{aligned}
\tag{8}
\]
where
\[
\begin{aligned}
c_{Z,4}
&:=
4L_f^a+6
+12(L_\sigma^a)^2
+6(L_\sigma^x)^2,\\
b_{Z,4}^x
&:=
(L_f^x)^4+6(L_\sigma^x)^2.
\end{aligned}
\tag{9}
\]

Since
\[
(1+r)^4\le 8(1+r^4),
\qquad r\ge0,
\]
we have
\[
\begin{aligned}
\mathbb E
\left(1+\|\bar a_t^{\theta,x}\|\right)^4
&\le
8\left(
1+\mathbb E\|\bar a_t^{\theta,x}\|^4
\right)\\
&\le
8(1+M_{A,4}).
\end{aligned}
\]
Consequently,
\[
\frac{d}{dt}\mathbb E\|Z_t\|^4
\le
c_{Z,4}\mathbb E\|Z_t\|^4
+
b_{Z,4}^x\|x'-x\|^4
+
b_{Z,4}^\theta\|\theta'-\theta\|^4,
\tag{10}
\]
where
\[
b_{Z,4}^\theta
:=
8B_{\nabla f}^4(1+M_{A,4}).
\tag{11}
\]

Since \(Z_0=0\), Gr\"onwall's inequality yields
\[
\mathbb E\|Z_t\|^4
\le
\frac{e^{c_{Z,4}t}-1}{c_{Z,4}}
\left[
b_{Z,4}^x\|x'-x\|^4
+
b_{Z,4}^\theta\|\theta'-\theta\|^4
\right].
\tag{12}
\]
Taking fourth roots and using
\[
(u+v)^{1/4}\le u^{1/4}+v^{1/4},
\qquad u,v\ge0,
\]
gives
\[
\begin{aligned}
\left(\mathbb E\|Z_t\|^4\right)^{1/4}
\le{}&
\left[
\frac{e^{c_{Z,4}t}-1}{c_{Z,4}}
b_{Z,4}^x
\right]^{1/4}
\|x'-x\|\\
&+
\left[
\frac{e^{c_{Z,4}t}-1}{c_{Z,4}}
b_{Z,4}^\theta
\right]^{1/4}
\|\theta'-\theta\|.
\end{aligned}
\tag{13}
\]

Therefore, setting
\[
C_{A,4}^x
:=
\left[
\frac{e^{c_{Z,4}T}-1}{c_{Z,4}}
b_{Z,4}^x
\right]^{1/4},
\qquad
C_{A,4}^\theta
:=
\left[
\frac{e^{c_{Z,4}T}-1}{c_{Z,4}}
b_{Z,4}^\theta
\right]^{1/4},
\tag{14}
\]
and taking the supremum over \(t\in[0,T]\) proves \((A4)\).
Evaluating the estimate at \(t=T\) proves \((A4')\).

\end{proof}
\end{lem}


% \begin{remark}
% In the exact notation of Appendix B, Lemma B.4 gives the corresponding projected estimate for \(f,g\in\mathcal F_{K,s}\):
% \[
% \left(
% \mathbb E\|\widetilde A_T^{f,x}-\widetilde A_T^{g,x'}\|^2
% \right)^{1/2}
% \lesssim
% (1+\ell_2)\sqrt{KT}e^{C_3KT}
% \left\{\|x-x'\|+d_{K,R}(f,g)\right\}.
% \]
% For a parameterized class, if
% \[
% d_{K,R}(f_\theta,f_{\theta'})
% \le
% L_{\Theta,R}\|\theta-\theta'\|,
% \]
% then the above gives \(C_A^x\) and \(C_A^\theta\), up to localization.
% \end{remark}
\begin{lem}[Score-integrand regularity]
Define
\[
G_\theta(t,a,x)
:=
\nabla_\theta f_\theta(t,a,x)\sigma(t,a,x)^{-1}.
\]
Claim
\[
\|G_\theta(t,a,x)\|_{\rm HS}\le M_G(1+\|a\|),
\]
where $M_G \coloneqq \kappa^{-1/2}B_{\nabla f}$
and
\[
\|G_{\theta'}(t,a',x')-G_\theta (t,a,x)\|_{\rm HS}\leq L_G^\theta\|\theta'-\theta\|
+
(L_G^{a,1}+L_G^{a,2}(1+\|a\|))\|a'-a\|
+
(L_G^{x,1}+L_G^{x,2}(1+\|a\|))\|x'-x\|.
\]
\begin{proof}
    For the boundedness, We have 
\[
\|G_\theta(t,a,x)\|_{\rm HS}\leq \|\nabla_\theta f_\theta(t,a,x)\|_{\rm HS}\|\sigma(t,a,x)^{-1}\|_{\rm op}\leq B_{\nabla f}(1+\|a\|) \kappa^{-\frac{1}{2}}
\]

For the L-lipschitzness, note that $\sigma'^{-1}-\sigma^{-1}=\sigma'^{-1}(\sigma-\sigma')\sigma^{-1}$, so $\|\sigma^{-1}(t,a',x')-\sigma^{-1}(t,a,x)\|_{\rm op}\leq\kappa^{-1}\|\sigma(t,a',x')-\sigma(t,a,x)\|_{\rm HS}\leq \kappa^{-1}(L^a_\sigma \|a'-a\|+L^x_\sigma\|x'-x\|)$
\begin{align}
    &\|G_{\theta'}(t,a',x')-G_\theta (t,a,x)\|_{\rm HS}\\
    =&\|[\nabla_\theta f_{\theta'}(t,a',x')-\nabla_\theta f_\theta(t,a,x)]\sigma^{-1}(t,a',x')+\nabla_\theta f_\theta(t,a,x)[\sigma^{-1}(t,a',x')-\sigma^{-1}(t,a,x)]\|_{\rm HS}\\
    % =&\|[\nabla_\theta f_{\theta'}(t,a',x')-\nabla_\theta f_\theta(t,a,x)]\sigma^{-1}(t,a',x')+\nabla_\theta f_\theta(t,a,x)[\sigma^{-1}(t,a',x')-\sigma^{-1}(t,a,x)]\|_{\rm HS}\\
    \leq&\|\nabla_\theta f_{\theta'}(t,a',x')-\nabla_\theta f_\theta(t,a,x)\|_{\rm HS}\|\sigma^{-1}(t,a',x')\|_{\rm op}+\|\nabla_\theta f_\theta(t,a,x)\|_{\rm HS}\|\sigma^{-1}(t,a',x')-\sigma^{-1}(t,a,x)\|_{\rm op} \\
    \leq &L_G^\theta\|\theta'-\theta\|
+
(L_G^{a,1}+L_G^{a,2}(1+\|a\|))\|a'-a\|
+
(L_G^{x,1}+L_G^{x,2}(1+\|a\|))\|x'-x\|.
\end{align}

where $L^\theta_G=\kappa^{-1/2}L^\theta_{\nabla f}$; $L^{a,1}_G=\kappa^{-1/2}L^a_{\nabla f}$;$L^{a,2}_G=\kappa^{-1}B_{\nabla f}L^a_\sigma$; $L^{x,1}_G=\kappa^{-1/2}L^x_{\nabla f}$; $L^{x,2}_G=\kappa^{-1}B_{\nabla f}L^x_\sigma$
\end{proof}
\end{lem}

% \begin{lem}[Score moment and score stability]
% The score satisfies
% \[
% \left(
% \mathbb E\|S_\theta^x\|^2
% \right)^{1/2}
% \le
% M_S,
% \qquad
% M_S:=\sqrt{2TM_G^2(1+M_{A,2})}.
% \tag{S0}
% \]
% Moreover,
% \[
% \left(
% \mathbb E
% \|S_{\theta'}^{x'}-S_\theta^x\|^2
% \right)^{1/2}
% \le
% C_S^x\|x'-x\|
% +
% C_S^\theta\|\theta'-\theta\|,
% \tag{S}
% \]
% where one may take
% \[
% C_S^x
% =
% \sqrt{T}(L_G^x+L_G^aC_A^x),
% \qquad
% C_S^\theta
% =
% \sqrt{T}(L_G^\theta+L_G^aC_A^\theta).
% \]

% \begin{proof}
% By Itô isometry,
% \[
% \mathbb E\|S_\theta^x\|^2
% =
% \int_0^T
% \mathbb E
% \|G_\theta(t,\bar a_t^{\theta,x},x)\|_{\rm HS}^2\,dt
% \le
% 2TM_G^2(1+\sup_{\theta,t,x}\E\|\bar{a}_t^{\theta,x}\|^2)=2TM_G^2(1+M_{A,2}).
% \]
% This proves \((S0)\).

% For stability, under the same-Brownian coupling,
% \[
% S_{\theta'}^{x'}-S_\theta^x
% =
% \int_0^T
% \left[
% G_{\theta'}(t,\bar a_t^{\theta',x'},x')
% -
% G_\theta(t,\bar a_t^{\theta,x},x)
% \right]dB_t.
% \]
% Again by Itô isometry,
% \[
% \mathbb E
% \|S_{\theta'}^{x'}-S_\theta^x\|^2
% =
% \int_0^T
% \mathbb E
% \left\|
% G_{\theta'}(t,\bar a_t^{\theta',x'},x')
% -
% G_\theta(t,\bar a_t^{\theta,x},x)
% \right\|_{\rm HS}^2\,dt.
% \]
% Using the proven property of \(G_\theta\) and the action-path stability estimate gives

% \begin{align*}
%     &\mathbb E
% \|S_{\theta'}^{x'}-S_\theta^x\|^2\\
% =&\int_0^T
% \mathbb E
% \left\|
% G_{\theta'}(t,\bar a_t^{\theta',x'},x')
% -
% G_\theta(t,\bar a_t^{\theta,x},x)
% \right\|_{\rm HS}^2\,dt\\
% \leq& \int^T_0\E\left[\left(L_G^\theta\|\theta'-\theta\|
% +
% (L_G^{a,1}+L_G^{a,2}(1+\|\bar a_t^{\theta,x}\|))\|\bar a_t^{\theta',x'}-\bar a_t^{\theta,x}\|
% +
% (L_G^{x,1}+L_G^{x,2}(1+\|\bar a_t^{\theta,x}\|))\|x'-x\|\right)^2\right]dt\\
% \leq& 3TL_G^\theta\|\theta'-\theta\|^2
% +
% 3T\sup_{t,\theta,x}\E\left[(L_G^{a,1}+L_G^{a,2}(1+\|\bar a_t^{\theta,x}\|))^2\|\bar a_t^{\theta',x'}-\bar a_t^{\theta,x}\|^2\right]
% \\
% +&
% 6T[(L_G^{x,1})^2+(L_G^{x,2})^2(1+\sup_{t,\theta,x}\E\|\bar a_t^{\theta,x}\|^2)]\|x'-x\|^2\\
% \leq& 3TL_G^\theta\|\theta'-\theta\|^2
% +
% 6T(L_G^{a,1})^2M_{A,2}+12T(L_G^{a,2})^2M_{A,2}+12T(L_G^{a,2})^2\E\left[\|\bar a_t^{\theta,x}\|^2\|\bar a_t^{\theta',x'}-\bar a_t^{\theta,x}\|^2\right]
% \\
% +&
% 6T[(L_G^{x,1})^2+(L_G^{x,2})^2(1+M_{A,2})]\|x'-x\|^2\\
% \leq& 3TL_G^\theta\|\theta'-\theta\|^2
% +
% 6T(L_G^{a,1})^2M_{A,2}+12T(L_G^{a,2})^2M_{A,2}+12T(L_G^{a,2})^2\left(\E\left[\|\bar a_t^{\theta,x}\|^4\right]\right)^{1/2}\left(\E\left[|\bar a_t^{\theta',x'}-\bar a_t^{\theta,x}\|^4\right]\right)^{1/2}
% \\
% +&
% 6T[(L_G^{x,1})^2+(L_G^{x,2})^2(1+M_{A,2})]\|x'-x\|^2\\
% \leq & 3TL_G^\theta\|\theta'-\theta\|^2
% +
% 6T(L_G^{a,1})^2M_{A,2}+12T(L_G^{a,2})^2M_{A,2}+12T(L_G^{a,2})^2(M_{A,4})^{1/2}\left(\E\left[|\bar a_t^{\theta',x'}-\bar a_t^{\theta,x}\|^4\right]\right)^{1/2}
% \\
% +&
% 6T[(L_G^{x,1})^2+(L_G^{x,2})^2(1+M_{A,2})]\|x'-x\|^2\\
% \leq & 3TL_G^\theta\|\theta'-\theta\|^2
% +
% 6T(L_G^{a,1})^2M_{A,2}+12T(L_G^{a,2})^2M_{A,2}+24T(L_G^{a,2})^2(M_{A,4})^{1/2}\left((C^x_{A,4})^2\|x'-x\|^2+(C^\theta_{A,4})^2\|\theta'-\theta\|^2\right)
% \\
% +&
% 6T[(L_G^{x,1})^2+(L_G^{x,2})^2(1+M_{A,2})]\|x'-x\|^2\\
% \end{align*}
% \[
% \left(
% \mathbb E
% \|S_{\theta'}^{x'}-S_\theta^x\|^2
% \right)^{1/2}
% \le
% \sqrt{T}
% \left[
% (L_G^x+L_G^aC_A^x)\|x'-x\|
% +
% (L_G^\theta+L_G^aC_A^\theta)\|\theta'-\theta\|
% \right].
% \]
% \end{proof}
% \end{lem}

\textcolor{purple}{Renyuan take a look; also flow matching}
\begin{lem}[Score moment and stability]
Define
\[
S_\theta^x
:=
\int_0^T
G_\theta(t,\bar a_t^{\theta,x},x)\,dB_t.
\]
Then
\[
\sup_{\theta,x}
\mathbb E\|S_\theta^x\|^2
\le
2TM_G^2(1+M_{A,2}).
\tag{S0}
\]

Moreover,
\[
\left(
\mathbb E
\|S_{\theta'}^{x'}-S_\theta^x\|^2
\right)^{1/2}
\le
C_S^x\|x'-x\|
+
C_S^\theta\|\theta'-\theta\|,
\tag{S1}
\]
where
\[
C_S^x:=\sqrt{T}\,C_G^x,
\qquad
C_S^\theta:=\sqrt{T}\,C_G^\theta,
\]
with
\[
\begin{aligned}
C_G^x
:={}&
L_G^{x,1}
+
m_2L_G^{x,2}
+
L_G^{a,1}C_A^x
+
m_4L_G^{a,2}C_{A,4}^x,\\
C_G^\theta
:={}&
L_G^\theta
+
L_G^{a,1}C_A^\theta
+
m_4L_G^{a,2}C_{A,4}^\theta,
\end{aligned}
\tag{S2}
\]
and
\[
m_2:=\sqrt{2(1+M_{A,2})},
\qquad
m_4:=\left[8(1+M_{A,4})\right]^{1/4}
,\qquad M_S\coloneqq\sqrt{T}M_Gm_2\]

\begin{proof}
By It\^o's isometry,
\[
\begin{aligned}
\mathbb E\|S_\theta^x\|^2
&=
\int_0^T
\mathbb E
\|G_\theta(t,\bar a_t^{\theta,x},x)\|_{\rm HS}^2\,dt\\
&\le
M_G^2
\int_0^T
\mathbb E
\left(1+\|\bar a_t^{\theta,x}\|\right)^2dt\\
&\le
2TM_G^2
\left(
1+
\sup_{\theta,t,x}
\mathbb E\|\bar a_t^{\theta,x}\|^2
\right)\\
&=
2TM_G^2(1+M_{A,2}),
\end{aligned}
\]
which proves \((S0)\).

For the stability estimate, set
\[
Z_t
:=
\bar a_t^{\theta',x'}-\bar a_t^{\theta,x}.
\]
Under the same-Brownian coupling,
\[
\begin{aligned}
S_{\theta'}^{x'}-S_\theta^x
=
\int_0^T
\Big[
&G_{\theta'}(t,\bar a_t^{\theta',x'},x')\\
&-
G_\theta(t,\bar a_t^{\theta,x},x)
\Big]\,dB_t.
\end{aligned}
\]
Therefore, It\^o's isometry gives
\[
\begin{aligned}
&
\mathbb E
\|S_{\theta'}^{x'}-S_\theta^x\|^2\\
&=
\int_0^T
\mathbb E
\left\|
G_{\theta'}(t,\bar a_t^{\theta',x'},x')
-
G_\theta(t,\bar a_t^{\theta,x},x)
\right\|_{\rm HS}^2dt.
\end{aligned}
\tag{1}
\]

Applying \((A)\) with
\[
a'=\bar a_t^{\theta',x'},
\qquad
a=\bar a_t^{\theta,x},
\]
and then using Minkowski's inequality in \(L^2\), we obtain
\[
\begin{aligned}
&
\left(
\mathbb E
\left\|
G_{\theta'}(t,\bar a_t^{\theta',x'},x')
-
G_\theta(t,\bar a_t^{\theta,x},x)
\right\|_{\rm HS}^2
\right)^{1/2}\\
&\le
L_G^\theta\|\theta'-\theta\|
+
L_G^{a,1}
\left(\mathbb E\|Z_t\|^2\right)^{1/2}\\
&\quad+
L_G^{a,2}
\left(
\mathbb E
\left[
(1+\|\bar a_t^{\theta,x}\|)^2
\|Z_t\|^2
\right]
\right)^{1/2}\\
&\quad+
\left[
L_G^{x,1}
+
L_G^{x,2}
\left(
\mathbb E
(1+\|\bar a_t^{\theta,x}\|)^2
\right)^{1/2}
\right]\|x'-x\|.
\end{aligned}
\tag{2}
\]

We have
\[
\begin{aligned}
\left(
\mathbb E
(1+\|\bar a_t^{\theta,x}\|)^2
\right)^{1/2}
&\le
\sqrt{
2\left(
1+\mathbb E\|\bar a_t^{\theta,x}\|^2
\right)
}\\
&\le
m_2.
\end{aligned}
\tag{3}
\]
Moreover
\[
\begin{aligned}
&
\left(
\mathbb E
\left[
(1+\|\bar a_t^{\theta,x}\|)^2
\|Z_t\|^2
\right]
\right)^{1/2}\\
&\le
\left(
\mathbb E
(1+\|\bar a_t^{\theta,x}\|)^4
\right)^{1/4}
\left(
\mathbb E\|Z_t\|^4
\right)^{1/4}\\
&\le
m_4
\left[
C_{A,4}^x\|x'-x\|
+
C_{A,4}^\theta\|\theta'-\theta\|
\right].
\end{aligned}
\tag{4}
\]
% Since
% \[
% (1+r)^4\le 8(1+r^4),
% \qquad r\ge0,
% \]
% we have
% \[
% \left(
% \mathbb E
% (1+\|\bar a_t^{\theta,x}\|)^4
% \right)^{1/4}
% \le
% m_4.
% \]
% Thus
% \[
% \begin{aligned}
% &
% \left(
% \mathbb E
% \left[
% (1+\|\bar a_t^{\theta,x}\|)^2
% \|Z_t\|^2
% \right]
% \right)^{1/2}\\

% \end{aligned}
% \tag{4}
% \]

Combining \((2)\)--\((4)\) with the \(L^2\) action-path stability
estimate \((A2)\), we obtain, uniformly in \(t\in[0,T]\),
\[
\begin{aligned}
&
\left(
\mathbb E
\left\|
G_{\theta'}(t,\bar a_t^{\theta',x'},x')
-
G_\theta(t,\bar a_t^{\theta,x},x)
\right\|_{\rm HS}^2
\right)^{1/2}\\
&\le
C_G^x\|x'-x\|
+
C_G^\theta\|\theta'-\theta\|.
\end{aligned}
\tag{5}
\]

Finally, substituting \((5)\) into \((1)\) gives
\[
\begin{aligned}
\left(
\mathbb E
\|S_{\theta'}^{x'}-S_\theta^x\|^2
\right)^{1/2}
&\le
\sqrt{T}
\left[
C_G^x\|x'-x\|
+
C_G^\theta\|\theta'-\theta\|
\right]\\
&=
C_S^x\|x'-x\|
+
C_S^\theta\|\theta'-\theta\|.
\end{aligned}
\]
This proves \((S1)\).
\end{proof}
\end{lem}

\subsection{Synchronous coupling convention}
\label{subsec:synchronous-mdp-coupling}

We now fix the common-noise construction used in the value-function and
discounted-occupancy arguments. Let
\[
\left\{(\zeta_j,B^j,U_j)\right\}_{j\ge 0}
\]
be independent across MDP time steps \(j\), where
\(\zeta_j\sim\mu_0\), \(B^j=(B_t^j)_{0\le t\le T}\) is a standard
Brownian motion, and \(U_j\sim{\rm Unif}[0,1]\). For each \(j\), assume
that \(\zeta_j\), \(B^j\), and \(U_j\) are mutually independent and
independent of the coupled initial states.

Fix \(\vartheta_1,\vartheta_2\in\Theta\) and an arbitrary coupling
\((X_0^1,X_0^2)\) of their initial-state distributions. Define
\[
\mathcal F_j
:=
\sigma\left(
X_0^1,X_0^2,
\left\{(\zeta_k,B^k,U_k):0\le k<j\right\}
\right).
\]
Then \(X_j^1,X_j^2\) are \(\mathcal F_j\)-measurable, whereas
\((\zeta_j,B^j,U_j)\) is independent of \(\mathcal F_j\).

At MDP time step \(j\), conditionally on \(\mathcal F_j\), define the two
within-step action diffusions using the same fresh initial action noise
\(\zeta_j\) and the same fresh Brownian motion \(B^j\):
\[
\begin{aligned}
d\bar a_{j,t}^{\,i}
&=
f_{\vartheta_i}
\bigl(t,\bar a_{j,t}^{\,i},X_j^i\bigr)\,dt
+
\sigma
\bigl(t,\bar a_{j,t}^{\,i},X_j^i\bigr)\,dB_t^j,
\\
\bar a_{j,0}^{\,i}
&=\zeta_j,
\qquad
A_j^i:=\bar a_{j,T}^{\,i},
\qquad i\in\{1,2\}.
\end{aligned}
\label{eq:mdp-action-coupling}
\]
The next states are generated using the same fresh transition seed:
\[
X_{j+1}^i
=
\Phi_P(X_j^i,A_j^i,U_j),
\qquad i\in\{1,2\}.
\label{eq:mdp-transition-coupling}
\]

Because the randomness is fresh at each MDP time step, each marginal has the
correct MDP law. More precisely,
\[
\operatorname{Law}(A_j^i\mid\mathcal F_j)
=
\pi_{\vartheta_i}(\cdot\mid X_j^i).
\]
Let
\[
\mathcal G_j
:=
\mathcal F_j\vee\sigma(\zeta_j,B^j).
\]
Then \(X_j^i,A_j^i\) are \(\mathcal G_j\)-measurable and \(U_j\) is
independent of \(\mathcal G_j\), so
\[
\operatorname{Law}(X_{j+1}^i\mid\mathcal G_j)
=
P(\cdot\mid X_j^i,A_j^i).
\]
Thus each marginal process \(\{(X_j^i,A_j^i)\}_{j\ge0}\) is an MDP
trajectory generated by \(\pi_{\vartheta_i}\).
All stability estimates below are derived in their corresponding proofs
by conditioning first on \(\mathcal F_j\) for the action pair and then
on \(\mathcal G_j\) for the transition pair.

\begin{lem}[Value stability in state]
Define
\[
\alpha:=L_P^x+L_P^aC_A^x.
\]
Assume
\[
\gamma\alpha<1.
\tag{C}
\]
Then, uniformly in \(\theta\),
\[
|V_\theta(x)-V_\theta(y)|
\le
L_V^x\|x-y\|,
\tag{Vx}
\]
where
\[
L_V^x
:=
\frac{L_r^x+L_r^aC_A^x}{1-\gamma\alpha}.
\]

\begin{proof}
Apply the construction in
Subsection~\ref{subsec:synchronous-mdp-coupling} with
\[
\vartheta_1=\vartheta_2=\theta,
\qquad
X_0^1=x,
\qquad
X_0^2=y,
\]
and write
\[
(X_j,A_j):=(X_j^1,A_j^1),
\qquad
(Y_j,\widetilde A_j):=(X_j^2,A_j^2).
\]

Let
\[
D_j:=\left(\mathbb E\|X_j-Y_j\|^2\right)^{1/2},
\qquad
E_j:=\left(\mathbb E\|A_j-\widetilde A_j\|^2\right)^{1/2}.
\]
Conditionally on \(\mathcal F_j\), the current states are fixed.
Applying action-path stability at \(t=T\) therefore gives
\[
\left(
\mathbb E\left[
\|A_j-\widetilde A_j\|^2
\mid\mathcal F_j
\right]
\right)^{1/2}
\le
C_A^x\|X_j-Y_j\|.
\]
Taking the outer \(L^2\)-norm gives
\[
E_j\le C_A^xD_j.
\]

Conditionally on \(\mathcal G_j\), the current states and actions are
fixed while \(U_j\) is fresh. The synchronous transition estimate gives
\[
\begin{aligned}
&\left(
\mathbb E\left[
\|X_{j+1}-Y_{j+1}\|^2
\mid\mathcal G_j
\right]
\right)^{1/2}
\\
&\qquad\le
L_P^x\|X_j-Y_j\|
+
L_P^a\|A_j-\widetilde A_j\|.
\end{aligned}
\]
Taking the outer \(L^2\)-norm and using Minkowski's inequality gives
\[
D_{j+1}
\le
L_P^xD_j+L_P^aE_j
\le
\alpha D_j.
\]
Since \(D_0=\|x-y\|\), induction gives
\[
D_j\le \alpha^j\|x-y\|,
\qquad
E_j\le C_A^x\alpha^j\|x-y\|.
\]
Therefore,
\[
\begin{aligned}
|V_\theta(x)-V_\theta(y)|
&\le
\sum_{j=0}^{\infty}
\gamma^j
\mathbb E
\left[
|r(X_j,A_j)-r(Y_j,\widetilde A_j)|
\right]
\\
&\le
\sum_{j=0}^{\infty}
\gamma^j
\left[
L_r^xD_j+L_r^aE_j
\right]
\\
&\le
(L_r^x+L_r^aC_A^x)
\sum_{j=0}^{\infty}
(\gamma\alpha)^j
\|x-y\|.
\end{aligned}
\]
Since \(\gamma\alpha<1\), this proves \((Vx)\).
\end{proof}
\end{lem}

\begin{lem}[Value stability in parameter]
Let
\[
\alpha:=L_P^x+L_P^aC_A^x,
\qquad
\beta:=L_P^aC_A^\theta.
\]
Assume \(\gamma\alpha<1\). Then
\[
\|V_{\theta'}-V_\theta\|_\infty
\le
L_V^\theta\|\theta'-\theta\|,
\tag{Vtheta}
\]
where one may take
\[
L_V^\theta
:=
\frac{L_r^aC_A^\theta}{1-\gamma}
+
\frac{(L_r^x+L_r^aC_A^x)\beta\gamma}
{(1-\gamma)(1-\gamma\alpha)}.
\]

\begin{proof}
Fix an initial state \(x\). Apply the construction in
Subsection~\ref{subsec:synchronous-mdp-coupling} with
\[
(\vartheta_1,\vartheta_2)=(\theta,\theta'),
\qquad
X_0^1=X_0^2=x.
\]
Write \((X_j^\theta,A_j^\theta):=(X_j^1,A_j^1)\) and
\((X_j^{\theta'},A_j^{\theta'}):=(X_j^2,A_j^2)\).

Define
\[
D_j:=
\left(\mathbb E\|X_j^{\theta'}-X_j^\theta\|^2\right)^{1/2},
\qquad
E_j:=
\left(\mathbb E\|A_j^{\theta'}-A_j^\theta\|^2\right)^{1/2}.
\]
Conditionally on \(\mathcal F_j\), action-path stability at \(t=T\)
gives
\[
\begin{aligned}
&\left(
\mathbb E\left[
\|A_j^{\theta'}-A_j^\theta\|^2
\mid\mathcal F_j
\right]
\right)^{1/2}
\\
&\qquad\le
C_A^x\|X_j^{\theta'}-X_j^\theta\|
+
C_A^\theta\|\theta'-\theta\|.
\end{aligned}
\]
Taking the outer \(L^2\)-norm and using Minkowski's inequality gives
\[
E_j
\le
C_A^xD_j+C_A^\theta\|\theta'-\theta\|.
\]

Conditionally on \(\mathcal G_j\), the synchronous transition estimate
gives
\[
\begin{aligned}
&\left(
\mathbb E\left[
\|X_{j+1}^{\theta'}-X_{j+1}^\theta\|^2
\mid\mathcal G_j
\right]
\right)^{1/2}
\\
&\qquad\le
L_P^x\|X_j^{\theta'}-X_j^\theta\|
+
L_P^a\|A_j^{\theta'}-A_j^\theta\|.
\end{aligned}
\]
Taking the outer \(L^2\)-norm gives
\[
D_{j+1}
\le
L_P^xD_j+L_P^aE_j
\le
\alpha D_j+\beta\|\theta'-\theta\|.
\]
Since \(D_0=0\), iteration yields
\[
D_j
\le
\beta
\sum_{k=0}^{j-1}\alpha^k
\|\theta'-\theta\|.
\]
Hence
\[
\begin{aligned}
|V_{\theta'}(x)-V_\theta(x)|
&\le
\sum_{j=0}^{\infty}
\gamma^j
\left[
L_r^xD_j+L_r^aE_j
\right]
\\
&\le
\sum_{j=0}^{\infty}
\gamma^j
\left[
(L_r^x+L_r^aC_A^x)D_j
+
L_r^aC_A^\theta\|\theta'-\theta\|
\right].
\end{aligned}
\]
Using
\[
\sum_{j=0}^{\infty}\gamma^j=\frac{1}{1-\gamma}
\]
and
\[
\sum_{j=0}^{\infty}
\gamma^j
\sum_{k=0}^{j-1}\alpha^k
=
\frac{\gamma}{(1-\gamma)(1-\gamma\alpha)},
\]
we obtain
\[
|V_{\theta'}(x)-V_\theta(x)|
\le
L_V^\theta\|\theta'-\theta\|.
\]

with \[
L_V^\theta
:=
\frac{L_r^aC_A^\theta}{1-\gamma}
+
\frac{(L_r^x+L_r^aC_A^x)\beta\gamma}
{(1-\gamma)(1-\gamma\alpha)}.
\]
Taking the supremum over \(x\) proves the result.
\end{proof}
\end{lem}


\begin{lem}[\(Q\)-function bounds]
The \(Q\)-function satisfies
\[
\|Q_\theta\|_\infty
\le
M_Q,
\qquad
M_Q:=\frac{r_{\max}}{1-\gamma}.
\tag{Q0}
\]
Moreover,
\[
|Q_\theta(x,a)-Q_\theta(x',a')|
\le
L_Q^x\|x-x'\|+L_Q^a\|a-a'\|,
\tag{Qx}
\]
where
\[
L_Q^x:=L_r^x+\gamma L_V^xL_P^x,
\qquad
L_Q^a:=L_r^a+\gamma L_V^xL_P^a.
\]
Finally,
\[
\|Q_{\theta'}-Q_\theta\|_\infty
\le
L_Q^\theta\|\theta'-\theta\|,
\tag{Qtheta}
\]
where
\[
L_Q^\theta:=\gamma L_V^\theta.
\]
Consequently,
\[
|Q_{\theta'}(x',a')-Q_\theta(x,a)|
\le
L_Q^x\|x'-x\|
+
L_Q^a\|a'-a\|
+
L_Q^\theta\|\theta'-\theta\|.
\tag{Qjoint}
\]

\begin{proof}
The uniform bound follows from the boundedness of the reward. For all
\((x,a)\in\mathcal X\times\mathcal A\),
\[
|Q_\theta(x,a)|
\le
\sum_{j=0}^{\infty}\gamma^j r_{\max}
=
\frac{r_{\max}}{1-\gamma}
=
M_Q.
\]

Next, fix \(x,x'\in\mathcal X\) and \(a,a'\in\mathcal A\). By the
Bellman representation,
\[
Q_\theta(x,a)
=
r(x,a)
+
\gamma
\int_{\mathcal X}
V_\theta(z)P(dz\mid x,a).
\]
Therefore,
\[
\begin{aligned}
|Q_\theta(x,a)-Q_\theta(x',a')|
&\le
|r(x,a)-r(x',a')|
\\
&\quad+
\gamma
\left|
\int_{\mathcal X}V_\theta(z)P(dz\mid x,a)
-
\int_{\mathcal X}V_\theta(z)P(dz\mid x',a')
\right|.
\end{aligned}
\]

Let \(U\sim{\rm Unif}[0,1]\), and define the synchronously coupled next
states
\[
Z:=\Phi_P(x,a,U),
\qquad
Z':=\Phi_P(x',a',U).
\]
Then
\[
Z\sim P(\cdot\mid x,a),
\qquad
Z'\sim P(\cdot\mid x',a'),
\]
and hence
\[
\begin{aligned}
&\left|
\int_{\mathcal X}V_\theta(z)P(dz\mid x,a)
-
\int_{\mathcal X}V_\theta(z)P(dz\mid x',a')
\right|
\\
&=
\left|
\mathbb E_U
\left[
V_\theta(Z)-V_\theta(Z')
\right]
\right|
\\
&\le
L_V^x\mathbb E_U\|Z-Z'\|
\\
&\le
L_V^x
\left(
\mathbb E_U\|Z-Z'\|^2
\right)^{1/2}
\\
&\le
L_V^x
\left(
L_P^x\|x-x'\|
+
L_P^a\|a-a'\|
\right).
\end{aligned}
\]
Together with reward Lipschitzness, this gives
\[
|Q_\theta(x,a)-Q_\theta(x',a')|
\le
L_Q^x\|x-x'\|
+
L_Q^a\|a-a'\|,
\]
where
\[
L_Q^x
=
L_r^x+\gamma L_V^xL_P^x,
\qquad
L_Q^a
=
L_r^a+\gamma L_V^xL_P^a.
\]

For parameter stability, fix the same state-action pair \((x,a)\).
Because the reward and transition kernel do not depend directly on
\(\theta\),
\[
\begin{aligned}
Q_{\theta'}(x,a)-Q_\theta(x,a)
&=
\gamma
\int_{\mathcal X}
\left[
V_{\theta'}(z)-V_\theta(z)
\right]
P(dz\mid x,a).
\end{aligned}
\]
Consequently,
\[
\begin{aligned}
|Q_{\theta'}(x,a)-Q_\theta(x,a)|
&\le
\gamma
\|V_{\theta'}-V_\theta\|_\infty
\\
&\le
\gamma L_V^\theta\|\theta'-\theta\|.
\end{aligned}
\]
Taking the supremum over \((x,a)\) proves
\[
\|Q_{\theta'}-Q_\theta\|_\infty
\le
L_Q^\theta\|\theta'-\theta\|,
\qquad
L_Q^\theta:=\gamma L_V^\theta.
\]

Finally,
\[
\begin{aligned}
|Q_{\theta'}(x',a')-Q_\theta(x,a)|
&\le
|Q_{\theta'}(x',a')-Q_{\theta'}(x,a)|
\\
&\quad+
|Q_{\theta'}(x,a)-Q_\theta(x,a)|
\\
&\le
L_Q^x\|x'-x\|
+
L_Q^a\|a'-a\|
+
L_Q^\theta\|\theta'-\theta\|.
\end{aligned}
\]
This proves \((Qjoint)\).
\end{proof}
\end{lem}

\begin{lem}[\(\Psi\)-Lipschitzness]
The map
\[
(\theta,x)\mapsto \Psi_\theta(x)
\]
is Lipschitz. More precisely,
\[
\|\Psi_{\theta'}(x')-\Psi_\theta(x)\|
\le
L_{\Psi,x}\|x'-x\|
+
L_{\Psi,\theta}\|\theta'-\theta\|,
\tag{Psi}
\]
where one may take
\[
L_{\Psi,x}
:=
M_QC_S^x
+
M_S(L_Q^x+L_Q^aC_A^x),
\]
and
\[
L_{\Psi,\theta}
:=
M_QC_S^\theta
+
M_S(L_Q^aC_A^\theta+L_Q^\theta).
\]

\begin{proof}
Use the same-Brownian coupling for the two state-conditioned action diffusions. Then
\[
\begin{aligned}
\Psi_{\theta'}(x')-\Psi_\theta(x)
&=
\mathbb E
\left[
S_{\theta'}^{x'}Q_{\theta'}(x',A_{\theta'}^{x'})
-
S_\theta^xQ_\theta(x,A_\theta^x)
\right]
\\
&=
\mathbb E
\left[
(S_{\theta'}^{x'}-S_\theta^x)Q_\theta(x,A_\theta^x)
\right]
\\
&\quad+
\mathbb E
\left[
S_{\theta'}^{x'}
\left(
Q_{\theta'}(x',A_{\theta'}^{x'})
-
Q_\theta(x,A_\theta^x)
\right)
\right].
\end{aligned}
\]
Therefore,
\[
\begin{aligned}
\|\Psi_{\theta'}(x')-\Psi_\theta(x)\|
&\le
M_Q
\left(
\mathbb E
\|S_{\theta'}^{x'}-S_\theta^x\|^2
\right)^{1/2}
\\
&\quad+
M_S
\left(
\mathbb E
\left|
Q_{\theta'}(x',A_{\theta'}^{x'})
-
Q_\theta(x,A_\theta^x)
\right|^2
\right)^{1/2}.
\end{aligned}
\]
Using score stability, action stability, and the joint \(Q\)-bound, we get
\[
\|\Psi_{\theta'}(x')-\Psi_\theta(x)\|
\le
\left[
M_QC_S^x
+
M_S(L_Q^x+L_Q^aC_A^x)
\right]\|x'-x\|
\]
\[
+
\left[
M_QC_S^\theta
+
M_S(L_Q^aC_A^\theta+L_Q^\theta)
\right]\|\theta'-\theta\|.
\]
This proves the claim.
\end{proof}
\end{lem}

\begin{lem}[Discounted occupancy stability]
\label{lem:discounted-occupancy-stability}
Let
\[
d_\rho^{\pi_\theta}(dx)
:=
(1-\gamma)\sum_{j=0}^{\infty}\gamma^j
\mathbb P_\rho^{\pi_\theta}(X_j^\theta\in dx).
\]
Set
\[
\alpha:=L_P^x+L_P^aC_A^x,
\qquad
\beta:=L_P^aC_A^\theta,
\qquad
C_d
:=
\frac{\gamma\beta}{1-\gamma\alpha}
=
\frac{\gamma L_P^aC_A^\theta}
{1-\gamma(L_P^x+L_P^aC_A^x)}.
\]
If
\[
\gamma\alpha<1,
\tag{4.7-c}
\]
then for any bounded vector-valued test function
\[
\varphi:\mathcal X\to\mathbb R^p
\]
satisfying
\[
\|\varphi(x)-\varphi(y)\|
\le
L_\varphi\|x-y\|,
\]
we have
\[
\left\|
\int_{\mathcal X}\varphi(x)
\left[
d_\rho^{\pi_{\theta'}}-d_\rho^{\pi_\theta}
\right](dx)
\right\|
\le
L_\varphi C_d\|\theta'-\theta\|,
\tag{4.7}
\]
% In particular,
% \[
% W_1(d_\rho^{\pi_{\theta'}},d_\rho^{\pi_\theta})
% \le
% C_d\|\theta'-\theta\|.
% \]


\begin{proof}
Apply the coupling convention from
Subsection~\ref{subsec:synchronous-mdp-coupling} with
\[
(\vartheta_1,\vartheta_2)=(\theta,\theta'),
\qquad
X_0^1=X_0^2\sim\rho
\quad\text{almost surely}.
\]
Write \((X_j^\theta,A_j^\theta):=(X_j^1,A_j^1)\) and
\((X_j^{\theta'},A_j^{\theta'}):=(X_j^2,A_j^2)\), and define
\[
D_j
:=
\left(
\mathbb E_{\rm sync}
\|X_j^{\theta'}-X_j^\theta\|^2
\right)^{1/2},
\qquad
E_j
:=
\left(
\mathbb E_{\rm sync}
\|A_j^{\theta'}-A_j^\theta\|^2
\right)^{1/2}.
\]

Conditionally on \(\mathcal F_j\), action-path stability at \(t=T\)
gives
\[
\begin{aligned}
&\left(
\mathbb E_{\rm sync}\left[
\|A_j^{\theta'}-A_j^\theta\|^2
\mid\mathcal F_j
\right]
\right)^{1/2}
\\
&\qquad\le
C_A^x\|X_j^{\theta'}-X_j^\theta\|
+
C_A^\theta\|\theta'-\theta\|.
\end{aligned}
\]
Taking the outer \(L^2\)-norm gives
\[
E_j
\le
C_A^xD_j+C_A^\theta\|\theta'-\theta\|.
\]

Conditionally on \(\mathcal G_j\), the synchronous transition estimate
gives
\[
\begin{aligned}
&\left(
\mathbb E_{\rm sync}\left[
\|X_{j+1}^{\theta'}-X_{j+1}^\theta\|^2
\mid\mathcal G_j
\right]
\right)^{1/2}
\\
&\qquad\le
L_P^x\|X_j^{\theta'}-X_j^\theta\|
+
L_P^a\|A_j^{\theta'}-A_j^\theta\|.
\end{aligned}
\]
Taking the outer \(L^2\)-norm and using the preceding action estimate
yields
\[
D_{j+1}
\le
L_P^xD_j+L_P^aE_j
\le
\alpha D_j+\beta\|\theta'-\theta\|.
\tag{4.8}
\]
Since \(D_0=0\), iteration gives
\[
D_j
\le
\beta
\sum_{k=0}^{j-1}\alpha^k
\|\theta'-\theta\|.
\tag{4.9}
\]

Now expand the discounted occupancy measures term by term. For any Lipschitz test function \(\varphi\),
\[
\begin{aligned}
&\int_{\mathcal X}\varphi(x)
\left[
d_\rho^{\pi_{\theta'}}-d_\rho^{\pi_\theta}
\right](dx)
\\
&=
(1-\gamma)
\sum_{j=0}^{\infty}
\gamma^j
\left[
\mathbb E_{\rm sync}\varphi(X_j^{\theta'})
-
\mathbb E_{\rm sync}\varphi(X_j^\theta)
\right]
\\
&=
(1-\gamma)
\sum_{j=0}^{\infty}
\gamma^j
\mathbb E_{\rm sync}
\left[
\varphi(X_j^{\theta'})
-
\varphi(X_j^\theta)
\right].
\end{aligned}
\]
Therefore,
\[
\begin{aligned}
&\left\|
\int_{\mathcal X}\varphi(x)
\left[
d_\rho^{\pi_{\theta'}}-d_\rho^{\pi_\theta}
\right](dx)
\right\|
\\
&\le
(1-\gamma)
\sum_{j=0}^{\infty}
\gamma^j
\mathbb E_{\rm sync}
\left[
\|\varphi(X_j^{\theta'})
-
\varphi(X_j^\theta)\|
\right]
\\
&\le
L_\varphi
(1-\gamma)
\sum_{j=0}^{\infty}
\gamma^j
\mathbb E_{\rm sync}
\|X_j^{\theta'}-X_j^\theta\|
\\
&\le
L_\varphi
(1-\gamma)
\sum_{j=0}^{\infty}
\gamma^jD_j.
\end{aligned}
\]
Using \((4.9)\),
\[
\begin{aligned}
(1-\gamma)
\sum_{j=0}^{\infty}
\gamma^jD_j
&\le
(1-\gamma)
\sum_{j=0}^{\infty}
\gamma^j
\beta
\sum_{k=0}^{j-1}\alpha^k
\|\theta'-\theta\|
\\
&=
\beta
(1-\gamma)
\sum_{j=0}^{\infty}
\gamma^j
\sum_{k=0}^{j-1}\alpha^k
\|\theta'-\theta\|.
\end{aligned}
\]
Since
\[
(1-\gamma)
\sum_{j=0}^{\infty}
\gamma^j
\sum_{k=0}^{j-1}\alpha^k
=
\frac{\gamma}{1-\gamma\alpha},
\]
provided \(\gamma\alpha<1\), we obtain
\[
(1-\gamma)
\sum_{j=0}^{\infty}
\gamma^jD_j
\le
\frac{\gamma\beta}{1-\gamma\alpha}
\|\theta'-\theta\|.
\]
Thus
\[
\left\|
\int_{\mathcal X}\varphi(x)
\left[
d_\rho^{\pi_{\theta'}}-d_\rho^{\pi_\theta}
\right](dx)
\right\|
\le
L_\varphi
\frac{\gamma\beta}{1-\gamma\alpha}
\|\theta'-\theta\|.
\]
This proves \((4.7)\).\\
\textcolor{purple}{Another way of proofing}
Let $P_\theta(x,dy) \coloneqq \int_\mathcal{A} P(dy|x,a)\pi_\theta(da|x) = \E[P(dy|x,a^{\theta,x}_T)]$. Then 
\[d^{\pi_\theta}_\rho=(1-\gamma)(I-\gamma P^*_\theta)^{-1}\rho\]
Define $R_\theta=\sum_{k=0}^\infty\gamma^k (P_\theta)^k=(I-\gamma P_\theta)^{-1}$ and $R^*_\theta=\sum_{k=0}^\infty\gamma^k (P^*_\theta)^k=(I-\gamma P^*_\theta)^{-1}$, we have 
\begin{align*}
    d^{\pi_{\theta'}}_\rho-d^{\pi_\theta}_\rho&=(1-\gamma)[R_{\theta'}^*-R^*_\theta]\rho\\
    &=(1-\gamma)R_{\theta'}^*\gamma[P^*_{\theta'}-P^*_\theta]R_\theta^*\rho\\
    &=\gamma R^*_{\theta'}[P^*_{\theta'}-P^*_\theta]d^{\pi_\theta}_\rho
\end{align*}
Then 
\[\int_\mathcal{X}\varphi(x)[d^{\pi_{\theta'}}_\rho-d^{\pi_\theta}_\rho](dx)=\gamma\int_\mathcal{X}[(P_{\theta'}-P_\theta)R_{\theta'}\varphi](x)d^{\pi_\theta}_\rho(dx)
\]
Note that $\text{Lip}(R_{\theta'}\varphi)\leq \sum^\infty_{k=0}\gamma^k\alpha^kL_{\varphi}=\frac{1}{1-\gamma \alpha}L_\varphi$
Then
\[
\|(P_{\theta'}-P_\theta)R_{\theta'}\varphi (x)\|_\infty\leq \frac{\beta}{1-\gamma \alpha}L_\varphi \|\theta'-\theta\|
\]

So
\[
\int_\mathcal{X}\varphi(x)[d^{\pi_{\theta'}}_\rho-d^{\pi_\theta}_\rho](dx)\leq  \frac{\gamma\beta}{1-\gamma \alpha}L_\varphi \|\theta'-\theta\|
\]
% Finally, taking the supremum over scalar \(1\)-Lipschitz test functions gives
% \[
% W_1(d_\rho^{\pi_{\theta'}},d_\rho^{\pi_\theta})
% \le
% C_d\|\theta'-\theta\|.
% \]


\end{proof}
\end{lem}

\subsection{Main theorem}

\begin{thm}[\(L\)-smoothness of \(J(\theta)=V^{\pi_\theta}(\rho)\)]

% \thm[\(L\)-smoothness of \(J(\theta)=V^{\pi_\theta}(\rho)\)]
Assume the preceding assumptions hold, and assume
\[
\gamma\alpha<1.
\tag{C}
\]
Then \(J\) is \(L\)-smooth:
\[
\|\nabla J(\theta')-\nabla J(\theta)\|
\le
L\|\theta'-\theta\|,
\]

% with
% \[
% L
% =
% \frac{L_{\Psi,\theta}+L_{\Psi,x}C_d}{1-\gamma}.
% \]

\begin{proof}
By the policy-gradient formula,
\[
\nabla J(\theta')
-
\nabla J(\theta)
=
\frac{1}{1-\gamma}
\left[
\int_{\mathcal X}\Psi_{\theta'}(x)d_\rho^{\pi_{\theta'}}(dx)
-
\int_{\mathcal X}\Psi_\theta(x)d_\rho^{\pi_\theta}(dx)
\right].
\]
Add and subtract
\[
\int_{\mathcal X}\Psi_{\theta'}(x)d_\rho^{\pi_\theta}(dx).
\]
Then
\[
\nabla J(\theta')-\nabla J(\theta)
=
\frac{1}{1-\gamma}(I+II),
\]
where
\[
I
:=
\int_{\mathcal X}
\left[
\Psi_{\theta'}(x)-\Psi_\theta(x)
\right]d_\rho^{\pi_\theta}(dx),
\]
and
\[
II
:=
\int_{\mathcal X}
\Psi_{\theta'}(x)
\left[
d_\rho^{\pi_{\theta'}}-d_\rho^{\pi_\theta}
\right](dx).
\]

First, since \(d_\rho^{\pi_\theta}\) is a probability measure and \(\Psi\) is Lipschitz in \(\theta\),
\[
\begin{aligned}
\|I\|
&\le
\int_{\mathcal X}
\|\Psi_{\theta'}(x)-\Psi_\theta(x)\|d_\rho^{\pi_\theta}(dx)
\\
&\le
L_{\Psi,\theta}\|\theta'-\theta\|.
\end{aligned}
\tag{1}
\]

% Next, let \(N\sim{\rm Geom}(1-\gamma)\), and use the synchronous coupling so that
% \[
% X_N^{\theta'}\sim d_\rho^{\pi_{\theta'}},
% \qquad
% X_N^\theta\sim d_\rho^{\pi_\theta}.
% \]
% Then
% \[
% II
% =
% \mathbb E_{\rm sync}
% \left[
% \Psi_{\theta'}(X_N^{\theta'})
% -
% \Psi_{\theta'}(X_N^\theta)
% \right].
% \]
% Using Lipschitzness of \(\Psi_{\theta'}\) in \(x\),
% \[
% \begin{aligned}
% \|II\|
% &\le
% \mathbb E_{\rm sync}
% \left[
% \|\Psi_{\theta'}(X_N^{\theta'})
% -
% \Psi_{\theta'}(X_N^\theta)\|
% \right]
% \\
% &\le
% L_{\Psi,x}
% \mathbb E_{\rm sync}
% \|X_N^{\theta'}-X_N^\theta\|.
% \end{aligned}
% \]
% By discounted occupancy stability,
% \[
% \|II\|
% \le
% L_{\Psi,x}C_d\|\theta'-\theta\|.
% \tag{2}
% \]
Applying Lemma~\ref{lem:discounted-occupancy-stability} with
\(\varphi=\Psi_{\theta'}\),
\[
\|II\|
\le
L_{\Psi,x}C_d\|\theta'-\theta\|.
\tag{2}
\]
Combining \((1)\) and \((2)\),
\[
\begin{aligned}
\|\nabla J(\theta')-\nabla J(\theta)\|
&\le
\frac{1}{1-\gamma}
\left(
\|I\|+\|II\|
\right)
\\
&\le
\frac{L_{\Psi,\theta}+L_{\Psi,x}C_d}{1-\gamma}
\|\theta'-\theta\|.
\end{aligned}
\]
Thus \(J\) is \(L\)-smooth with
\[
L
=
\frac{L_{\Psi,\theta}+L_{\Psi,x}C_d}{1-\gamma}.
\]
\end{proof}
\end{thm}
\section{Sample-based analysis}
% \begin{lem}[Unbiased gradient sampling]
%     Define $J_H(\theta)\coloneqq \E[\sum^{H-1}_{t=0}\gamma^tr(X_t,A_t)]$ and 
%     \\$V^\theta_{k,H}(x)\coloneqq \E_\theta[\sum^{H-1}_{t=k} \gamma^{t-k} r(X_t,A_t)|X_k=x]$.
%     Define $\hat{g}_H=\sum^{H-1}_{k=1}S_{k,\theta}G_{k,H}$, where $S_{k,\theta}=\int^T_0 \nabla_\theta f_\theta(t,\bar{a}^{\theta,X_k}_{k,t})dB_t^k$ and $G_{k,H}=\sum^{H-1}_{k=0}\gamma^tr(X_t,A_t)$, we claim
%     \[
%     \nabla_\theta J_H(\theta)=\E[\hat{g}_H]
%     \]
%     \begin{proof}
%     Using a performance difference lemma for $Q^\theta_{k,H}(x,a)=r(x,a)+\gamma \int_\mathcal{X}V^{\theta}_{k+1,H}P(dy|x,a)$, 
%     \[
%     J_H(\theta+\epsilon \beta)-J_H(\theta)=\sum^{H-1}_{k=0} \gamma^k \E_{\theta+\epsilon \beta} [Q^\theta_{k,H}(X_k,a)[\pi_{\theta+\epsilon\beta}-\pi_\theta](da|X_k)]
%     \]
%     Follows a similar calculation as in the paper, we have 

%     \begin{align}
%        \nabla J_H(\theta)&=\sum^{H-1}_{k=0} \gamma^k \E_\theta[S_{k,\theta}Q^\theta_{k,H}(X_k,A_k)]\\
%        &=\E[\sum^{H-1}_{k=0}\gamma^kS_{k,\theta}\E[G_{k,H}|X_k,A_k]]\\
%        &=\E[\hat{g_H}]
%     \end{align}
%     % \begin{align}
%     %     \E[\hat{g}_H]&=\E[\sum^{H-1}_{k=1}S_{k,\theta}G_{k,H}]\\
%     %     &=\\
%     %     &=\E[\sum^{H-1}_{k=1}S_{k,\theta}Q^{\pi_\theta}(X_k,A_k)]\\
%     %     &=\nabla_\theta J_H(\theta)
%     % \end{align}
        
%     \end{proof}
% \end{lem}

\begin{lem}[Unbiased truncated policy-gradient estimator]
\label{lem:unbiased-truncated-gradient}

Fix \(H\in\mathbb N\), and define the truncated objective
\[
J_H(\theta)
:=
\mathbb E_\theta
\left[
\sum_{t=0}^{H-1}
\gamma^t r(X_t,A_t)
\right].
\]

For \(0\le k\le H\), define
\[
V_{k,H}^{\theta}(x)
:=
\mathbb E_\theta
\left[
\sum_{t=k}^{H-1}
\gamma^{t-k}r(X_t,A_t)
\,\middle|\,
X_k=x
\right],
\qquad
V_{H,H}^{\theta}\equiv 0,
\]
and, for \(0\le k\le H-1\),
\[
Q_{k,H}^{\theta}(x,a)
:=
r(x,a)
+
\gamma
\int_{\mathcal X}
V_{k+1,H}^{\theta}(y)
P(dy\mid x,a).
\]

At environment time \(k\), conditionally on \(X_k\), let
\[
d\bar a_{k,s}^{\theta,X_k}
=
f_\theta
\bigl(s,\bar a_{k,s}^{\theta,X_k},X_k\bigr)\,ds
+
\sigma
\bigl(s,\bar a_{k,s}^{\theta,X_k},X_k\bigr)\,dB_s^k,
\]
and set
\[
A_k
:=
\bar a_{k,T}^{\theta,X_k}.
\]

Define the corresponding Girsanov score
\[
S_{k,\theta}
:=
\int_0^T
\nabla_\theta f_\theta
\bigl(s,\bar a_{k,s}^{\theta,X_k},X_k\bigr)
\sigma
\bigl(s,\bar a_{k,s}^{\theta,X_k},X_k\bigr)^{-1}
\,dB_s^k,
\]
and the truncated return-to-go
\[
G_{k,H}
:=
\sum_{t=k}^{H-1}
\gamma^{t-k}r(X_t,A_t).
\]

% Assume that the hypotheses of the one-action Girsanov derivative
% formula hold and that, for every \(k,\theta,x\),
% the map
% \[
% a\longmapsto Q_{k,H}^{\theta}(x,a)
% \]
% is bounded and continuous. Define
% \[
% \widehat g_H(\theta)
% :=
% \sum_{k=0}^{H-1}
% \gamma^kS_{k,\theta}G_{k,H}.
% \]
Finally, define the estimator
\[
\hat{g}_H\coloneqq\sum^{H-1}_{k=0}\gamma^k S_{k,\theta}G_{k,H}
\]
Then
\[
\mathbb E_\theta[\widehat g_H(\theta)]
=
\nabla_\theta J_H(\theta).
\]
In particular, for 
\[
\hat{g}_{m,H}\coloneqq \frac{1}{m}\sum^{m}_{i=1}\sum^{H-1}_{k=0}\gamma^k S^{(i)}_{k,\theta}G^{(i)}_{k,H}
\]
we have 
\[
\E_\theta[\hat{g}_{m,H}]=\nabla_\theta J_H(\theta)
\]
\end{lem}

\begin{proof}
Fix a direction \(\beta\in\mathbb R^p\), and set
\[
\theta_\varepsilon
:=
\theta+\varepsilon\beta.
\]

By the finite-horizon performance-difference identity,
\[
\begin{aligned}
J_H(\theta+\epsilon \beta)-J_H(\theta)
=
\sum_{k=0}^{H-1}
\gamma^k
\mathbb E_{\theta_\varepsilon}
\left[
\int_{\mathcal A}
Q_{k,H}^{\theta}(X_k,a)
\bigl(
\pi_{\theta_\varepsilon}
-
\pi_\theta
\bigr)(da\mid X_k)
\right].
\end{aligned}
\tag{1}
\]

% Here \(Q_{k,H}^{\theta}\) is frozen at the base parameter
% \(\theta\). For fixed \(k\) and \(x\), define
% \[
% F_{k,\varepsilon}(x)
% :=
% \frac{1}{\varepsilon}
% \int_{\mathcal A}
% Q_{k,H}^{\theta}(x,a)
% \bigl(
% \pi_{\theta_\varepsilon}
% -
% \pi_\theta
% \bigr)(da\mid x).
% \]

% Applying the one-action Girsanov derivative formula with the fixed
% test function
% \[
% \phi(a)=Q_{k,H}^{\theta}(x,a)
% \]
% gives
% \[
% \begin{aligned}
% \lim_{\varepsilon\to0}
% F_{k,\varepsilon}(x)
% =
% \mathbb E_\theta
% \left[
% \left\langle
% S_\theta^x,\beta
% \right\rangle
% Q_{k,H}^{\theta}(x,A_\theta^x)
% \right].
% \end{aligned}
% \tag{2}
% \]

% Using the finite-horizon state-law stability and the uniform
% integrability supplied by the boundedness of
% \(Q_{k,H}^{\theta}\) and the score moment bounds, we may pass to the
% limit in (1). Hence
% \[
% \begin{aligned}
% \left\langle
% \nabla_\theta J_H(\theta),\beta
% \right\rangle
% =
% \sum_{k=0}^{H-1}
% \gamma^k
% \mathbb E_\theta
% \left[
% \left\langle
% S_{k,\theta},\beta
% \right\rangle
% Q_{k,H}^{\theta}(X_k,A_k)
% \right].
% \end{aligned}
% \]

% Since this holds for every \(\beta\in\mathbb R^p\),

% Using a performance difference lemma for $Q^\theta_{k,H}(x,a)=r(x,a)+\gamma \int_\mathcal{X}V^{\theta}_{k+1,H}P(dy|x,a)$, 
% \[
% J_H(\theta+\epsilon \beta)-J_H(\theta)=\sum^{H-1}_{k=0} \gamma^k \E_{\theta+\epsilon \beta} [Q^\theta_{k,H}(X_k,a)[\pi_{\theta+\epsilon\beta}-\pi_\theta](da|X_k)]
% \]
Follows a similar calculation as in the paper, we have 
% \begin{align*}
%    \nabla J_H(\theta)&=\sum^{H-1}_{k=0} \gamma^k \E_\theta[S_{k,\theta}Q^\theta_{k,H}(X_k,A_k)]\\
%    &=\E[\sum^{H-1}_{k=0}\gamma^kS_{k,\theta}\E[G_{k,H}|X_k,A_k]]\\
%    &=\E[\hat{g_H}]
% \end{align*}
\[
\nabla_\theta J_H(\theta)
=
\sum_{k=0}^{H-1}
\gamma^k
\mathbb E_\theta
\left[
S_{k,\theta}
Q_{k,H}^{\theta}(X_k,A_k)
\right].
\tag{3}
\]

Let \(\mathcal G_k\) be the sigma-field generated by the MDP history
up to \(X_k\) and by the entire internal action-diffusion path used
at environment time \(k\). Then
\[
X_k,\qquad A_k,\qquad S_{k,\theta}
\]
are \(\mathcal G_k\)-measurable. By the Markov property and the
freshness of the environment and future policy randomness,
\[
\mathbb E_\theta
\left[
G_{k,H}
\,\middle|\,
\mathcal G_k
\right]
=
Q_{k,H}^{\theta}(X_k,A_k).
\tag{4}
\]

Therefore,
\[
\begin{aligned}
\mathbb E_\theta
\left[
S_{k,\theta}G_{k,H}
\right]
&=
\mathbb E_\theta
\left[
S_{k,\theta}
\mathbb E_\theta
\left[
G_{k,H}
\,\middle|\,
\mathcal G_k
\right]
\right]
\\
&=
\mathbb E_\theta
\left[
S_{k,\theta}
Q_{k,H}^{\theta}(X_k,A_k)
\right].
\end{aligned}
\tag{5}
\]

Substituting (5) into (3) yields
\[
\begin{aligned}
\nabla_\theta J_H(\theta)
&=
\mathbb E_\theta
\left[
\sum_{k=0}^{H-1}
\gamma^kS_{k,\theta}G_{k,H}
\right]
\\
&=
\mathbb E_\theta[\widehat g_H(\theta)].
\end{aligned}
\]

This proves the claim.
\end{proof}
\begin{lem}[Variance of the estimator] 
\label{lem:truncated-gradient-variance}

% Suppose that
% \[
% \mathbb E_\theta
% \left[
% S_{k,\theta}
% \,\middle|\,
% \mathcal F_k
% \right]
% =0
% \]
% and
% \[
% \mathbb E_\theta
% \left[
% \|S_{k,\theta}\|^2
% \,\middle|\,
% \mathcal F_k
% \right]
% \le M_S^2
% \]
% almost surely, uniformly over \(k\) and \(\theta\), where
% \(\mathcal F_k\) is the information available immediately before the
% action diffusion at environment time \(k\) is sampled. 

% Recall the one-trajectory estimator
% \[
% \widehat g_H(\theta)
% :=
% \sum_{k=0}^{H-1}
% \gamma^k S_{k,\theta}G_{k,H},
% \qquad
% G_{k,H}
% :=
% \sum_{t=k}^{H-1}
% \gamma^{t-k}r(X_t,A_t),
% \]
% and suppose that
% \[
% \mathbb E_\theta[\widehat g_H(\theta)]
% =
% \nabla_\theta J_H(\theta).
% \]

% Then
We have
\[
\mathbb E_\theta
\left[
\left\|
\widehat g_H(\theta)
-
\nabla_\theta J_H(\theta)
\right\|^2
\right]
\le
\frac{r_{\max}^2M_S^2}{(1-\gamma)^3}.
\tag{V}
\]

Moreover, if
\[
\widehat g_{m,H}(\theta)
:=
\frac{1}{m}
\sum_{i=1}^m
\widehat g_H^{(i)}(\theta),
\]
where the \(m\) trajectories are independent, then
\[
\mathbb E_\theta
\left[
\left\|
\widehat g_{m,H}(\theta)
-
\nabla_\theta J_H(\theta)
\right\|^2
\right]
\le
\frac{r_{\max}^2M_S^2}
{m(1-\gamma)^3}.
\tag{V-m}
\]
with $M_S\coloneqq\sqrt{T}M_Gm_2$

\begin{proof}

If \(i<j\), then \(S_{i,\theta}\) is
\(\mathcal F_j\)-measurable. Therefore,
\[
\begin{aligned}
\mathbb E_\theta
\left[
\left\langle
S_{i,\theta},S_{j,\theta}
\right\rangle
\right]
&=
\mathbb E_\theta
\left[
\left\langle
S_{i,\theta},
\mathbb E_\theta
\left[
S_{j,\theta}
\,\middle|\,
\mathcal F_j
\right]
\right\rangle
\right]
\\
&=0.
\end{aligned}
\]

Hence,
\[
\begin{aligned}
\mathbb E_\theta
\left[
\|\sum_{k=0}^t S_{k,\theta}\|^2
\right]
&=
\sum_{k=0}^t
\mathbb E_\theta
\left[
\|S_{k,\theta}\|^2
\right]
\\
&\le
(t+1)M_S^2.
\end{aligned}
\tag{1}
\]

Note that
\[
\begin{aligned}
\widehat g_H(\theta)
&=
\sum_{k=0}^{H-1}
S_{k,\theta}
\sum_{t=k}^{H-1}
\gamma^t r(X_t,A_t)
\\
&=
\sum_{t=0}^{H-1}
\gamma^t r(X_t,A_t)\sum_{k=0}^t S_{k,\theta}.
\end{aligned}
\tag{2}
\]

By Cauchy--Schwarz inequality,
\[
\begin{aligned}
\|\widehat g_H(\theta)\|^2
&\le
\left(
\sum_{t=0}^{H-1}\gamma^t
\right)
\left(
\sum_{t=0}^{H-1}
\gamma^t
r(X_t,A_t)^2
\|\sum_{k=0}^t S_{k,\theta}\|^2
\right).
\end{aligned}
\]

Taking expectations and using \((1)\),
\[
\begin{aligned}
\mathbb E_\theta
\left[
\|\widehat g_H(\theta)\|^2
\right]
&\le
\frac{r_{\max}^2(1-\gamma^H)}{1-\gamma}
\sum_{t=0}^{H-1}
\gamma^t
\mathbb E_\theta
\left[
\|\sum_{k=0}^t S_{k,\theta}\|^2
\right]
\\
&\le
\frac{r_{\max}^2M_S^2(1-\gamma^H)}{1-\gamma}
\sum_{t=0}^{H-1}
(t+1)\gamma^t
\\
&\le
\frac{r_{\max}^2M_S^2(1-\gamma^H)}{(1-\gamma)^3}\left(1-\gamma^H-H\gamma^H(1-\gamma)\right)
\\
&\le
\frac{r_{\max}^2M_S^2}{(1-\gamma)^3},
\end{aligned}
\tag{3}
\]
% where we used
% \[
% \sum_{t=0}^{\infty}(t+1)\gamma^t
% =
% \frac{1}{(1-\gamma)^2}.
% \]

Since
\[
\mathbb E_\theta[\widehat g_H(\theta)]
=
\nabla_\theta J_H(\theta),
\]
we have
\[
\begin{aligned}
\mathbb E_\theta
\left[
\left\|
\widehat g_H(\theta)
-
\nabla_\theta J_H(\theta)
\right\|^2
\right]
&\le
\mathbb E_\theta
\left[
\|\widehat g_H(\theta)\|^2
\right]
\\
&\le
\frac{r_{\max}^2M_S^2}{(1-\gamma)^3}.
\end{aligned}
\]

% For the mini-batch estimator, set
% \[
% Y_i
% :=
% \widehat g_H^{(i)}(\theta)
% -
% \nabla_\theta J_H(\theta).
% \]
% The \(Y_i\)'s are independent and centered, so
Then
\[
\begin{aligned}
\mathbb E_\theta
\left[
\left\|
\widehat g_{m,H}(\theta)
-
\nabla_\theta J_H(\theta)
\right\|^2
\right]
&=
\frac{1}{m^2}
\sum_{i=1}^m
\mathbb E_\theta[\|\widehat g_H^{(i)}(\theta)
-
\nabla_\theta J_H(\theta)\|^2]
\\
&\le
\frac{r_{\max}^2M_S^2}
{m(1-\gamma)^3}.
\end{aligned}
\]

This completes the proof.
\end{proof}    
\end{lem}
\begin{lem}[Truncation Error]
Claim that
\[
\sup_\theta\|\nabla J(\theta)-\nabla J_H(\theta)\|^2\leq B_H
\]
with $B_H=\frac{\gamma^{2H}r_{max}^2 M_S^2}{(1-\gamma)^2}\left(H+1+\frac{\gamma}{1-\gamma}\right)$
\begin{proof}
    \begin{align*}
        \|\nabla J(\theta)-\nabla J_H(\theta)\|^2&=\left\|\E\left[\sum^\infty_{t=H} \gamma^t\sum^t_{k=0}r_tS_{k,\theta}\right]\right\|^2\\
        &\leq\E\left[\left\|\sum^\infty_{t=H} \gamma^t\sum^t_{k=0}r_tS_{k,\theta}\right\|^2\right]\\
        &\leq \frac{r^2_{max}\gamma^H}{1-\gamma}\E\left[\sum^\infty_{t=H} \gamma^t\left\|\sum^t_{k=0}S_{k,\theta}\right\|^2\right]\\
        &\leq \frac{\gamma^{2H}r_{max}^2 M_S^2}{(1-\gamma)^2}\left(H+1+\frac{\gamma}{1-\gamma}\right)
    \end{align*}
\end{proof}
\end{lem}
\begin{lem}[Stationarity]
    Consider the stochastic gradient ascent
    \[
    \theta_{n+1}=\theta_n+\eta \hat{g}_{m,H}(\theta_n)
    \]
    with $\eta\leq 1/L$, we claim 
    \begin{equation*}
        \frac{1}{N}\sum^{N-1}_{n=0}\E\|\nabla J(\theta_n)\|^2\leq \frac{2\sup_{\theta\in\Theta} (J(\theta)- J(\theta_0))}{\eta N}+B_H+\frac{L\eta}{m}\frac{r_{\max}^2M_S^2}
{(1-\gamma)^3}
    \end{equation*}
    Particularly at $\eta=\frac{1}{L}$, we have    
    \begin{equation*}
        \frac{1}{N}\sum^{N-1}_{n=0}\E\|\nabla J(\theta_n)\|^2\leq \frac{2L\sup_{\theta\in\Theta} (J(\theta)- J(\theta_0))}{ N}+B_H+\frac{1}{m}\frac{r_{\max}^2M_S^2}
{(1-\gamma)^3}
    \end{equation*}
    
\end{lem}
\begin{proof}
By the L-smoothness,
\begin{equation*}
    J(\theta_{n+1})\geq J(\theta_n)+\langle \eta \hat{g}_{m,H}(\theta_n),\nabla J(\theta_n)\rangle -\frac{L\eta^2}{2}\| \hat{g}_{m,H}\|^2
\end{equation*}
Taking conditional expectation on both side, we have
\begin{equation*}
    \E[J(\theta_{n+1})|\theta_n]\geq J(\theta_n)+\eta\langle  \nabla J_H(\theta_n),\nabla J(\theta_n)\rangle -\frac{L\eta^2}{2}\left(\|\nabla J_H(\theta_n)\|^2+\frac{r_{\max}^2M_S^2}
{m(1-\gamma)^3}\right)
\end{equation*}
and then with $\eta\leq 1/L$
\begin{equation*}
    \|\nabla J(\theta_n)\|^2\leq \frac{2}{\eta}(\E[J(\theta_{n+1})|\theta_n]-J(\theta_n))+B_H+\frac{L\eta}{m}\frac{r_{\max}^2M_S^2}
{(1-\gamma)^3}
\end{equation*}
Telescoping and we get the claimed result
\end{proof}

% \begin{lem}[Projected update Stationarity]
%         Consider the stochastic gradient ascent
%     \[
%     \theta_{n+1}=\Pi_\Theta(\theta_n+\eta \hat{g}_{m,H}(\theta_n))
%     \]
% \end{lem}
\newpage


\subsection{Stationarity of projected stochastic gradient ascent}

Let $\Theta\subseteq\mathbb R^p$ be nonempty, closed, and convex,
and let $\Pi_\Theta$ denote its Euclidean projection. Define the
\emph{population gradient mapping} for maximization by
\begin{equation}
\mathcal G_\eta(\theta)
:=\frac{\Pi_\Theta(\theta+\eta\nabla J(\theta))-\theta}{\eta},
\qquad \theta\in\Theta,\quad \eta>0.
\label{eq:proj-pg-gradient-mapping}
\end{equation}

\begin{lem}[One-step projected ascent with gradient error]
\label{lem:proj-pg-one-step}
% Suppose $J$ is differentiable on a neighborhood of $\Theta$ and
% its gradient is $L$-Lipschitz on $\Theta$, with $L>0$.
For any $\theta\in\Theta$, $\widehat g\in\mathbb R^p$, and
$0<\eta\le L^{-1}$, set
\[
\theta^+:=\Pi_\Theta(\theta+\eta\widehat g),
\qquad e:=\widehat g-\nabla J(\theta).
\]
Then
\begin{equation}
J(\theta^+)-J(\theta)
\ge \frac{\eta}{2}\|\mathcal G_\eta(\theta)\|^2
     -\frac{\eta}{2}\|e\|^2.
\label{eq:proj-pg-one-step}
\end{equation}
\end{lem}

\begin{proof}
Write $g:=\nabla J(\theta)$ and introduce the exact-gradient
projected point
\[
\theta^\star:=\Pi_\Theta(\theta+\eta g).
\]
Define, for $v\in\Theta$,
\[
q(v):=\langle g,v-\theta\rangle
       -\frac{\|v-\theta\|^2}{2\eta},
\qquad
\widehat q(v):=q(v)+\langle e,v-\theta\rangle.
\]
The point $\theta^+$ maximizes the $\eta^{-1}$-strongly concave
function $\widehat q$ over $\Theta$. The projection property gives
\[
\widehat q(\theta^+)
\ge \widehat q(\theta^\star)
     +\frac{\|\theta^+-\theta^\star\|^2}{2\eta}.
\]
Consequently, completing the square yields
\begin{align*}
q(\theta^+)
&\ge q(\theta^\star)
   +\frac{\|\theta^+-\theta^\star\|^2}{2\eta}
   -\langle e,\theta^+-\theta^\star\rangle\\
&\ge q(\theta^\star)-\frac{\eta}{2}\|e\|^2.
\end{align*}

The projection variational inequality for $\theta^\star$, tested
against the feasible point $\theta$, gives
\[
\langle \theta+\eta g-\theta^\star,
        \theta-\theta^\star\rangle\le0.
\]
Hence
\[
\langle g,\theta^\star-\theta\rangle
\ge\frac{\|\theta^\star-\theta\|^2}{\eta},
\qquad
q(\theta^\star)
\ge\frac{\|\theta^\star-\theta\|^2}{2\eta}
=\frac{\eta}{2}\|\mathcal G_\eta(\theta)\|^2.
\]
Finally, $L$-smoothness along the feasible line segment and
$\eta\le L^{-1}$ imply
\begin{align*}
J(\theta^+)-J(\theta)
&\ge\langle g,\theta^+-\theta\rangle
     -\frac L2\|\theta^+-\theta\|^2\\
&\ge q(\theta^+).
\end{align*}
Combining these inequalities proves
\eqref{eq:proj-pg-one-step}.
\end{proof}

\begin{thm}[Expected stationarity of projected stochastic ascent]
\label{thm:proj-pg-stationarity}
Assume the smoothness hypotheses of Lemma~\ref{lem:proj-pg-one-step}
and suppose
\[
J_\Theta^\star:=\sup_{\theta\in\Theta}J(\theta)<\infty.
\]
Fix a deterministic $\theta_0\in\Theta$, integers $N,m,H\ge1$,
and a constant step size $0<\eta\le L^{-1}$.
% Let $\mathscr H_n$ be the information available before the batch at
% iteration $n$ is sampled, and
let
\[
\widehat g_n:=\widehat g_{m,H}^{(n)}(\theta_n)
\]
be formed from $m$ fresh, conditionally independent on-policy
% trajectories under $\pi_{\theta_n}$. Assume that the previously
% established conditional variance and gradient-truncation bounds hold:
% \begin{align}
% \mathbb E[\widehat g_n\mid\mathscr H_n]
% &=\nabla J_H(\theta_n),
% \label{eq:proj-pg-unbiased}\\
% \mathbb E\!\left[
%  \|\widehat g_n-\nabla J_H(\theta_n)\|^2
%  \mid\mathscr H_n\right]
% &\le\frac{\mathcal V}{m},
% \label{eq:proj-pg-variance}\\
% \sup_{\theta\in\Theta}
% \|\nabla J_H(\theta)-\nabla J(\theta)\|^2
% &\le B_H.
% \label{eq:proj-pg-bias}
% \end{align}
Consider the projected updates
\begin{equation}
\theta_{n+1}=\Pi_\Theta(\theta_n+\eta\widehat g_n),
\qquad n=0,\ldots,N-1.
\label{eq:proj-pg-update}
\end{equation}
Then, with $\Delta_J:=J_\Theta^\star-J(\theta_0)$,
\begin{equation}
\frac1N\sum_{n=0}^{N-1}
\mathbb E\|\mathcal G_\eta(\theta_n)\|^2
\le \frac{2\Delta_J}{\eta N}+B_H+\frac{1}{m}\frac{r_{\max}^2M_S^2}
{(1-\gamma)^3}.
\label{eq:proj-pg-stationarity-bound}
\end{equation}
% Equivalently, if $R$ is uniform on $\{0,\ldots,N-1\}$ and independent
% of all trajectory randomness, then
% \[
% \mathbb E\|\mathcal G_\eta(\theta_R)\|^2
% \le\frac{2\Delta_J}{\eta N}+B_H+\frac{1}{m}\frac{r_{\max}^2M_S^2}
% {(1-\gamma)^3}.
% \]
\end{thm}

\begin{proof}
Projection ensures $\theta_n\in\Theta$ almost surely for every $n$.
Define
\[
b_n:=\nabla J_H(\theta_n)-\nabla J(\theta_n),
\qquad
\xi_n:=\widehat g_n-\nabla J_H(\theta_n).
\]
Since $b_n$ is $\mathcal{F}_n$-measurable and
$\mathbb E[\xi_n\mid\mathcal{F}_n]=0$, the conditional cross term
vanishes. Thus
\begin{align*}
\mathbb E\!\left[
 \|\widehat g_n-\nabla J(\theta_n)\|^2
 \mid\mathcal{F}_n\right]
&=\|b_n\|^2+
  \mathbb E[\|\xi_n\|^2\mid\mathcal{F}_n]\\
&\le B_H+\frac{1}{m}\frac{r_{\max}^2M_S^2}
{(1-\gamma)^3}.
\end{align*}
Apply Lemma~\ref{lem:proj-pg-one-step} and take conditional
expectations to obtain
\[
\mathbb E[J(\theta_{n+1})\mid\mathcal{F}_n]
\ge J(\theta_n)
  +\frac{\eta}{2}\|\mathcal G_\eta(\theta_n)\|^2
  -\frac{\eta}{2}\left(B_H+\frac{1}{m}\frac{r_{\max}^2M_S^2}
{(1-\gamma)^3}\right).
\]
Taking total expectations and summing over $n=0,\ldots,N-1$ gives
\[
\frac{\eta}{2}\sum_{n=0}^{N-1}
\mathbb E\|\mathcal G_\eta(\theta_n)\|^2
\le \mathbb E[J(\theta_N)]-J(\theta_0)
  +\frac{N\eta}{2}\left(B_H+\frac{1}{m}\frac{r_{\max}^2M_S^2}
{(1-\gamma)^3}\right).
\]
Using $J(\theta_N)\le J_\Theta^\star$ and dividing by $N\eta/2$
proves \eqref{eq:proj-pg-stationarity-bound}. 
% The random-iterate
% statement follows from
% \[
% \mathbb E\|\mathcal G_\eta(\theta_R)\|^2
% =\frac1N\sum_{n=0}^{N-1}
%   \mathbb E\|\mathcal G_\eta(\theta_n)\|^2.
% \]
\end{proof}

% \begin{remark}[Stationarity and comparison with the unprojected result]
% For every $\eta>0$ and $\theta\in\Theta$,
% \[
% \mathcal G_\eta(\theta)=0
% \quad\Longleftrightarrow\quad
% \langle\nabla J(\theta),v-\theta\rangle\le0
% \quad\text{for all }v\in\Theta.
% \]
% Thus the conclusion concerns constrained first-order stationarity,
% not necessarily a vanishing ordinary gradient or global optimality.
% The mapping in \eqref{eq:proj-pg-gradient-mapping} uses the population
% gradient, not the sampled displacement
% $(\theta_{n+1}-\theta_n)/\eta$.

% At $\eta=L^{-1}$, the bound becomes
% \[
% \mathbb E\|\mathcal G_{1/L}(\theta_R)\|^2
% \le\frac{2L\Delta_J}{N}+B_H+\frac{\mathcal V}{m}.
% \]
% For smaller $\eta$, the present projected bound retains
% $\mathcal V/m$; it does not recover the sharper unprojected noise
% term $L\eta\mathcal V/m$. This is a limitation of this bound,
% not a claim that reducing the step size cannot help.

% In the exact-gradient case, $\widehat g_n=\nabla J(\theta_n)$,
% the error terms vanish and
% \[
% \frac1N\sum_{n=0}^{N-1}\|\mathcal G_\eta(\theta_n)\|^2
% \le\frac{2\Delta_J}{\eta N}.
% \]
% \end{remark}

% \subsubsection{Projected stationarity via stochastic composite optimization}
% \label{subsubsec:projected-stationarity-cited}

% We retain the standing assumptions and constants of
% Theorem~\ref{thm:J-smoothness},
% Lemmas~\ref{lem:unbiased-truncated-gradient},
% \ref{lem:truncated-gradient-variance}, and
% \ref{lem:gradient-truncation}. In particular, the condition
% $\gamma\alpha<1$ remains in force, and
% \[
%  L=\frac{L_{\Psi,\theta}+L_{\Psi,x}C_d}{1-\gamma}>0,
%  \qquad
%  \mathcal V=\frac{r_{\max}^2M_S^2}{(1-\gamma)^3},
%  \qquad
%  B_H=\mathcal V\gamma^{2H}\bigl(1+(1-\gamma)H\bigr).
% \]
% Here $L$ may be any positive upper bound supplied by the smoothness
% analysis. We take $\Theta\subset\mathbb R^p$ to be nonempty, closed, and
% convex. Closedness ensures that Euclidean projection is well defined;
% it imposes no additional regularity condition on the action SDE.

% The population smoothness and gradient-oracle estimates below are
% consequences of the preceding results, not additional assumptions about
% individual sampled objectives. To apply an unbiased-oracle theorem
% literally, we first record the finite-horizon specialization of the
% existing smoothness argument.

% \begin{lem}[Uniform finite-horizon smoothness]
% \label{lem:truncated-objective-smoothness}
% Under the standing hypotheses, for every $H\ge1$ and every
% $\theta,\theta'\in\Theta$,
% \begin{equation}
%  \|\nabla J_H(\theta')-\nabla J_H(\theta)\|
%  \le L\|\theta'-\theta\|.
%  \label{eq:uniform-truncated-smoothness}
% \end{equation}
% The constant $L$ is independent of $H$ and can be the same constant as
% in Theorem~\ref{thm:J-smoothness}.
% \end{lem}

% \begin{proof}
% For $0\le k<H$, use the finite-horizon value functions from
% Lemma~\ref{lem:unbiased-truncated-gradient} and define
% \[
%  \Psi_{k,H}^{\theta}(x)
%  :=\mathbb E\!\left[S_\theta^x
%              Q_{k,H}^{\theta}(x,A_\theta^x)\right].
% \]
% The preceding value-stability proofs remain valid with their geometric
% sums stopped at the remaining horizon: every upper-bounding summand is
% nonnegative. Hence the constants $L_V^x,L_V^\theta$ still bound
% $V_{k,H}^{\theta}$ uniformly in $k,H$. The Bellman definition of
% $Q_{k,H}^{\theta}$ then gives the same constants
% $M_Q,L_Q^x,L_Q^a,L_Q^\theta$ as in the preceding $Q$-function bounds.
% Repeating the score--$Q$ product estimate therefore yields
% \begin{equation}
%  \|\Psi_{k,H}^{\theta'}(x')-\Psi_{k,H}^{\theta}(x)\|
%  \le L_{\Psi,\theta}\|\theta'-\theta\|
%       +L_{\Psi,x}\|x'-x\|,
%  \label{eq:finite-horizon-Psi-stability}
% \end{equation}
% with no dependence on $k$ or $H$.

% The finite-horizon policy-gradient identity gives
% \[
%  \nabla J_H(\theta)
%  =\sum_{k=0}^{H-1}\gamma^k
%        \mathbb E\!\left[\Psi_{k,H}^{\theta}(X_k^\theta)\right].
% \]
% Under the same-initial-state synchronous coupling already used in the
% occupancy proof, set $\beta=L_P^aC_A^\theta$. Its state estimate is
% \[
%  \mathbb E\|X_k^{\theta'}-X_k^\theta\|
%  \le \beta\sum_{j=0}^{k-1}\alpha^j\|\theta'-\theta\|,
% \]
% where the inner sum is zero when $k=0$. Consequently,
% \begin{align*}
%  \|\nabla J_H(\theta')-\nabla J_H(\theta)\|
%  &\le \left[
%      L_{\Psi,\theta}\sum_{k=0}^{H-1}\gamma^k
%      +L_{\Psi,x}\beta\sum_{k=0}^{H-1}\gamma^k
%                          \sum_{j=0}^{k-1}\alpha^j
%                   \right]\|\theta'-\theta\|\\
%  &\le \left[
%        \frac{L_{\Psi,\theta}}{1-\gamma}
%        +\frac{L_{\Psi,x}\beta\gamma}
%                    {(1-\gamma)(1-\gamma\alpha)}
%                   \right]\|\theta'-\theta\|\\
%  &=L\|\theta'-\theta\|.
% \end{align*}
% The last equality uses $C_d=\gamma\beta/(1-\gamma\alpha)$.
% \end{proof}

% \begin{prop}[Projected stochastic policy-gradient stationarity]
% \label{prop:projected-stochastic-pg-stationarity}
% Work under the preceding standing hypotheses. Fix integers $N,m,H\ge1$
% and a deterministic initial parameter $\theta_0\in\Theta$. At every
% iteration, use the fresh on-policy batch in
% \eqref{eq:stationarity-batch}, and update
% \begin{equation}
%  \theta_{n+1}
%  =\Pi_\Theta\!\left(\theta_n+\eta\widehat g_{m,H}^{(n)}(\theta_n)\right),
%  \qquad n=0,\ldots,N-1,
%  \qquad \eta=\frac1{2L}.
%  \label{eq:projected-pg-update-cited}
% \end{equation}
% Define the two population gradient mappings
% \begin{equation}
%  \mathcal G_\eta^{H}(\theta)
%  :=\frac{\Pi_\Theta(\theta+\eta\nabla J_H(\theta))-\theta}{\eta},
%  \qquad
%  \mathcal G_\eta^{J}(\theta)
%  :=\frac{\Pi_\Theta(\theta+\eta\nabla J(\theta))-\theta}{\eta},
%  \label{eq:projected-population-mappings-cited}
% \end{equation}
% and set
% \[
%  \Delta_{J,H}:=\sup_{\vartheta\in\Theta}J_H(\vartheta)-J_H(\theta_0).
% \]
% Let $R$ be uniform on $\{0,\ldots,N-1\}$, independently of all trajectory
% randomness. Then
% \begin{equation}
%  \mathbb E\|\mathcal G_\eta^{H}(\theta_R)\|^2
%  =\frac1N\sum_{n=0}^{N-1}
%            \mathbb E\|\mathcal G_\eta^{H}(\theta_n)\|^2
%  \le \underbrace{\frac{8L\Delta_{J,H}}N
%                          +\frac{6\mathcal V}{m}}_{=:A_{N,m,H}}.
%  \label{eq:projected-truncated-GLZ-bound}
% \end{equation}
% Moreover, for the original infinite-horizon objective,
% \begin{align}
%  \mathbb E\|\mathcal G_\eta^{J}(\theta_R)\|^2
%  &=\frac1N\sum_{n=0}^{N-1}
%            \mathbb E\|\mathcal G_\eta^{J}(\theta_n)\|^2
%  \notag\\
%  &\le \left(\sqrt{A_{N,m,H}}+\sqrt{B_H}\right)^2
%  \notag\\
%  &\le \frac{16L\Delta_{J,H}}N
%                    +\frac{12\mathcal V}{m}+2B_H.
%  \label{eq:projected-stationarity-bound}
% \end{align}
% The supremum defining $\Delta_{J,H}$ need not be attained. Bounded
% rewards give
% \begin{equation}
%  0\le\Delta_{J,H}
%  \le\frac{2r_{\max}(1-\gamma^H)}{1-\gamma}
%  \le\frac{2r_{\max}}{1-\gamma}.
%  \label{eq:projected-gap-upper-bound}
% \end{equation}
% \end{prop}

% \begin{proof}
% Conditioning on $\mathscr H_n$, the parameter $\theta_n$ is fixed and
% the batch has the on-policy law used in the preceding estimator lemmas.
% Thus Lemma~\ref{lem:conditional-gradient-error} supplies, almost surely,
% \[
%  \mathbb E[\widehat g_n\mid\mathscr H_n]=\nabla J_H(\theta_n),
%  \qquad
%  \mathbb E\!\left[
%   \|\widehat g_n-\nabla J_H(\theta_n)\|^2
%   \mid\mathscr H_n\right]\le\frac{\mathcal V}{m}.
% \]

% Apply \citet[Corollary~3, Eq.~(44), in the arXiv version]
% {GhadimiLanZhang2016} to the minimization problem
% \[
%  \min_{\theta\in\Theta} f(\theta),\qquad f=-J_H,
% \]
% with zero nonsmooth objective term and Euclidean distance generator
% $\omega(\theta)=\|\theta\|^2/2$. Its strong-convexity modulus is one.
% A single-trajectory stochastic oracle is
% $-\widehat g_H(\theta)$; the cited oracle variance constant can be set
% to $\sigma_{\mathrm{oracle}}^2=\mathcal V$, and the batch size is $m$.
% Lemma~\ref{lem:truncated-objective-smoothness} verifies the required
% $L$-smoothness, and \eqref{eq:projected-gap-upper-bound} verifies the
% finite objective gap. In the cited notation,
% $D_\Psi^2=\Delta_{J,H}/L$. The constant step
% $\eta=1/(2L)$ makes the cited output distribution uniform.
% Its population projected gradient is $-\mathcal G_\eta^{H}$, whose
% norm is unchanged. Substitution into Eq.~(44) proves
% \eqref{eq:projected-truncated-GLZ-bound}.

% Next, nonexpansiveness of Euclidean projection (equivalently,
% \citealp[Proposition~1]{GhadimiLanZhang2016}) and
% Lemma~\ref{lem:gradient-truncation} imply, pointwise on $\Theta$,
% \begin{align*}
%  \|\mathcal G_\eta^{J}(\theta)-\mathcal G_\eta^{H}(\theta)\|
%  &\le \|\nabla J(\theta)-\nabla J_H(\theta)\|\\
%  &\le\sqrt{B_H}.
% \end{align*}
% The triangle inequality in $L^2$ gives
% \[
%  \left(\mathbb E\|\mathcal G_\eta^{J}(\theta_R)\|^2\right)^{1/2}
%  \le
%  \left(\mathbb E\|\mathcal G_\eta^{H}(\theta_R)\|^2\right)^{1/2}
%  +\sqrt{B_H}.
% \]
% This proves the first bound in \eqref{eq:projected-stationarity-bound};
% using $(a+b)^2\le2a^2+2b^2$ proves the separated bound. The stated
% averages equal the random-output expectations by independence and
% uniformity of $R$.
% \end{proof}

% \begin{prop}[A priori projected $\varepsilon$-stationarity]
% \label{prop:projected-stationarity-complexity-cited}
% Under Proposition~\ref{prop:projected-stochastic-pg-stationarity}, fix
% $\varepsilon>0$ and retain $\eta=1/(2L)$. It is sufficient to choose
% positive integers satisfying
% \begin{align}
%  N&\ge\max\!\left\{1,
%        \left\lceil\frac{96Lr_{\max}}{(1-\gamma)\varepsilon^2}
%                         \right\rceil\right\},
%  &
%  m&\ge\max\!\left\{1,
%        \left\lceil\frac{36\mathcal V}{\varepsilon^2}\right\rceil
%                        \right\},
%  \label{eq:projected-N-m-cited}\\
%  H&\ge\max\!\left\{1,
%        \left\lceil
%        \frac{\log(\max\{1,6\mathcal V/\varepsilon^2\})}
%             {\log(1/\gamma)}
%        \right\rceil\right\}.
%  \label{eq:projected-H-cited}
% \end{align}
% Then
% \[
%  \mathbb E\|\mathcal G_{1/(2L)}^{J}(\theta_R)\|^2
%  \le\varepsilon^2.
% \]
% If an upper bound on $\Delta_{J,H}$ is available, the first prescription
% may instead use $N\ge\max\{1,\lceil48L\Delta_{J,H}/\varepsilon^2\rceil\}$.
% \end{prop}

% \begin{proof}
% The first two prescriptions and
% \eqref{eq:projected-gap-upper-bound} make the optimization and sampling
% terms in \eqref{eq:projected-stationarity-bound} at most
% $\varepsilon^2/3$ each. Since
% \[
%  1+(1-\gamma)H\le e^{(1-\gamma)H}\le\gamma^{-H},
% \]
% we have $B_H\le\mathcal V\gamma^H$.
% The horizon prescription gives $B_H\le\varepsilon^2/6$, including the
% case $\mathcal V=0$, so the truncation term is also at most
% $\varepsilon^2/3$.
% \end{proof}
\newpage
% \section{Other route toards $S_1$}
% \subsection{Proving for piecewise affine function as multi-layer ReLU}
% \paragraph{Scope and notation.}


% Fix $T>0$, $K>0$, $\mathcal X=[0,1]^m$, a finite ReLU architecture with
% parameter vector $\theta\in\R^p$, and a bounded convex parameter set
% $\Theta$. All representations and coefficient maps below are defined on an
% open neighborhood of $\Theta$; the stated uniform bounds are required on
% $\Theta$. Let $y=(t,a,x)$ and $\widetilde y=(t,a^\top,x^\top,1)^\top$.

% \paragraph{Affine-region representation.}
% Assume a finite, consistently labelled collection $\mathscr S$ of branch
% templates such that
% \begin{equation}
%  \NN_\theta(t,a,x)
%  =M^1_{\theta,S}t+M^2_{\theta,S}a+M^3_{\theta,S}x+b_{\theta,S}
%  =\mathsf C_{\theta,S}\widetilde y,
%  \qquad S=\iota_\theta(t,a,x),
%  \label{eq:affine}
% \end{equation}
% where $\mathsf C_{\theta,S}=[M^1_{\theta,S}\ M^2_{\theta,S}\ M^3_{\theta,S}\ b_{\theta,S}]$.
% The branch labels are fixed, but their active regions may depend on $\theta$
% and may be empty. Affine templates are extended to all inputs, including
% points outside their active regions. Assume $\NN_\theta$ is continuous.

% \paragraph{Assumption A: drift and volatility.}
% Set $f_\theta=K\NN_\theta$. Assume constants $B_f,L_f^a,L_f^x,L_f^\theta$
% such that
% \begin{align}
%  \|f_\theta(t,a,x)\|&\le B_f(1+\|a\|),\\
%  \|f_\theta(t,a',x')-f_\theta(t,a,x)\|
%  &\le L_f^a\|a'-a\|+L_f^x\|x'-x\|,\\
%  \|f_{\theta'}(t,a,x)-f_\theta(t,a,x)\|
%  &\le L_f^\theta(1+\|a\|)\|\theta'-\theta\|.
%  \label{eq:drift-stability}
% \end{align}
% For the original class, the first two inequalities follow from its slope
% and origin bounds; the last is the fixed-architecture parameter perturbation
% estimate. Let $\sigma$ be measurable, symmetric, independent of $\theta$,
% and satisfy
% \begin{align}
%  \kappa I_d&\preceq\sigma(t,a,x)\sigma(t,a,x)^\top\preceq\Lambda I_d,
%  \qquad 0<\kappa\le\Lambda<\infty,\\
%  \|\sigma(t,a',x')-\sigma(t,a,x)\|_{\HS}
%  &\le L_\sigma^a\|a'-a\|+L_\sigma^x\|x'-x\|.
% \end{align}
% Let $\zeta$ be independent of the Brownian motion $B$, with
% $\E\|\zeta\|^4<\infty$. All comparisons use the same $(\zeta,B)$ in
% \begin{equation}
%  dA_t^{\theta,x}=f_\theta(t,A_t^{\theta,x},x)\,dt
%  +\sigma(t,A_t^{\theta,x},x)\,dB_t,
%  \qquad A_0^{\theta,x}=\zeta.
%  \label{eq:sde}
% \end{equation}

% \paragraph{Assumption B: regularity of affine templates.}
% Assume each $\mathsf C_{\theta,S}$ is $C^{1,1}$ in $\theta$, with
% \begin{align}
%  \sup_{\theta,S}\bigl(\|\mathsf C_{\theta,S}\|
%        +\|D_\theta\mathsf C_{\theta,S}\|\bigr)&\le B_C,\\
%  \|D_\theta\mathsf C_{\theta',S}
%        -D_\theta\mathsf C_{\theta,S}\|
%  &\le L_C\|\theta'-\theta\|.
%  \label{eq:branch-regularity}
% \end{align}
% Tensor norms are the Euclidean norms of the corresponding coefficient
% arrays. The bounds include templates that are inactive at one of the
% parameters. In a fixed ReLU architecture, activation-labelled affine
% coefficients are polynomial functions of the weights, so this assumption
% is available on bounded parameter boxes using an appropriate refinement.

% \paragraph{Assumption C: stable action-transverse switching description.}
% Suppose there are finitely many consistently indexed affine functions
% \begin{equation}
%  \ell_{\theta,j}(t,a,x)
%  =n_{\theta,j}^\top a+\alpha_{\theta,j}t
%        +\beta_{\theta,j}^\top x+c_{\theta,j},\qquad j=1,\ldots,J,
%  \label{eq:interfaces}
% \end{equation}
% with unit action normals $\|n_{\theta,j}\|=1$ and coefficient vectors
% $q_{\theta,j}=(n_{\theta,j},\alpha_{\theta,j},\beta_{\theta,j},c_{\theta,j})$
% satisfying
% \begin{equation}
%  \|q_{\theta,j}\|\le B_\Gamma,\qquad
%  \|q_{\theta',j}-q_{\theta,j}\|
%        \le L_\Gamma\|\theta'-\theta\|.
%  \label{eq:interface-regularity}
% \end{equation}
% There is a single map $\mathfrak r:\{0,1\}^J\to\mathscr S$, independent of
% $\theta$, such that off the switching hyperplanes
% \begin{equation}
%  \iota_\theta(t,a,x)
%  =\mathfrak r\bigl(\1_{\{\ell_{\theta,1}(t,a,x)>0\}},\ldots,
%             \1_{\{\ell_{\theta,J}(t,a,x)>0\}}\bigr).
%  \label{eq:selection}
% \end{equation}
% The hyperplanes may refine the actual region boundaries; not every zero
% set must be an actual interface. Globally fixed-sign tests can be omitted.
% The normalization is accompanied by the coefficient regularity in
% \eqref{eq:interface-regularity}; it is not a way to avoid a degeneracy
% assumption. State-only or time-only switching surfaces are not covered by
% this sufficient hypothesis.

% Define $H_\theta=\nabla_\theta f_\theta\in\R^{p\times d}$ off the switching
% set, using the branch formula, and assign any measurable representative
% of at most linear growth on that set. Let
% \begin{equation}
%  G_t^{\theta,x}=H_\theta(t,A_t^{\theta,x},x)
%              \sigma(t,A_t^{\theta,x},x)^{-1},\qquad
%  S_\theta^x=\int_0^T G_t^{\theta,x}\,dB_t.
% \end{equation}

% \begin{prop}[A7-av in the affine-region formulation]
% Under Assumptions A--C, there is a finite constant $C$, independent of the
% comparison pair, such that
% \begin{equation}
%  \int_0^T\E\|G_t^{\theta',x'}-G_t^{\theta,x}\|_{HS}\,dt
%  \le C\bigl(\|\theta'-\theta\|+\|x'-x\|\bigr)
%  \label{eq:a7av}
% \end{equation}
% for all $\theta,\theta'\in\Theta$ and $x,x'\in\mathcal X$.
% The constant may depend on the fixed architecture, $J,K,T,d,m,p$, all
% coefficient and interface bounds, the volatility constants, and
% $\E\|\zeta\|^4$. No architecture-uniform or degeneration-uniform bound is
% asserted.
% \end{prop}

% \begin{proof}
% Throughout, $C$ is a finite constant with only the stated dependencies.
% Fix a comparison pair and abbreviate
% $A_t=A_t^{\theta,x}$, $A_t'=A_t^{\theta',x'}$, $Z_t=A_t'-A_t$, and
% $\delta=\|\theta'-\theta\|+\|x'-x\|$. If $\delta=0$, pathwise uniqueness
% makes the conclusion immediate. Hence assume $\delta>0$.

% \paragraph{Step 1: maximal moments and a common random envelope.}
% Linear growth, bounded volatility, BDG, and Gronwall give
% \begin{equation}
%  \sup_{\theta,x}\E\sup_{t\le T}\|A_t^{\theta,x}\|^4\le C.
%  \label{eq:moment}
% \end{equation}
% The drift and diffusion differences in the equation for $Z$ satisfy
% \begin{align}
%  \|\Delta f_t\|&\le L_f^a\|Z_t\|+C\delta(1+\|A_t\|),\\
%  \|\Delta\sigma_t\|_{\HS}&\le L_\sigma^a\|Z_t\|+L_\sigma^x\delta.
% \end{align}
% For $h(u)=\E\sup_{t\le u}\|Z_t\|^4$, the integral equation, Holder, and BDG
% therefore imply
% \[
%  h(u)\le C\int_0^u h(v)\,dv
%        +C\delta^4\int_0^u(1+\E\|A_v\|^4)\,dv.
% \]
% Gronwall and \eqref{eq:moment} give
% \begin{equation}
%  \E\sup_{t\le T}\|Z_t\|^4\le C\delta^4.
%  \label{eq:path-stability}
% \end{equation}
% Consequently,
% \begin{equation}
%  R=1+\sup_{t\le T}\|A_t\|
%        +\delta^{-1}\sup_{t\le T}\|Z_t\|,
%  \qquad \E R^4\le C.
%  \label{eq:envelope}
% \end{equation}

% \paragraph{Step 2: deterministic-width occupation estimates.}
% Fix an interface index $j$ and define the globally affine projection
% \[
%  Y_t=\ell_{\theta,j}(t,A_t,x).
% \]
% This is an auxiliary process even where its hyperplane is not an active
% facet. Its dynamics are
% \[
%  dY_t=\mu_t\,dt+\nu_t\,dB_t,
%  \quad \mu_t=n_{\theta,j}^\top f_\theta(t,A_t,x)+\alpha_{\theta,j},
%  \quad \nu_t=n_{\theta,j}^\top\sigma(t,A_t,x).
% \]
% Thus $\kappa\,dt\le d\langle Y\rangle_t\le\Lambda\,dt$ and
% $\sup_t\|\mu_t\|_{L^4}\le C$. With local time normalized by
% \[
%  \int_0^T\varphi(Y_t)\,d\langle Y\rangle_t
%  =\int_\R\varphi(z)L_T^z(Y)\,dz,
% \]
% Tanaka's formula gives
% \[
%  L_T^z(Y)=|Y_T-z|-|Y_0-z|
%    -\int_0^T\operatorname{sgn}(Y_t-z)\mu_t\,dt
%    -\int_0^T\operatorname{sgn}(Y_t-z)\nu_t\,dB_t.
% \]
% The endpoint difference is bounded by $|Y_T-Y_0|$. Minkowski and BDG yield
% \begin{align*}
%  \|L_T^z(Y)\|_{L^4}
%  &\le 2\int_0^T\|\mu_t\|_{L^4}\,dt
%       +2C_{\rm BDG,4}\left\|\left(\int_0^T\|\nu_t\|^2dt\right)^{1/2}\right\|_{L^4}
%  \le C.
% \end{align*}
% This is uniform in $j,\theta,x,z$. Hence, for
% $O_j(r)=\int_0^T\ind_{\{|Y_t|\le r\}}dt$ and deterministic $r\ge0$,
% \begin{equation}
%  \|O_j(r)\|_{L^4}
%  \le\kappa^{-1}\int_{-r}^r\|L_T^z(Y)\|_{L^4}\,dz
%  \le Cr.
%  \label{eq:strip}
% \end{equation}
% The occupation formula also gives
% $\int_0^T\ind_{\{Y_t=0\}}dt=0$ almost surely. Since there are finitely many
% interfaces, the choice of derivative representative on their union does
% not affect $G$ as a $dt\,d\mathbb P$ equivalence class or its Ito integral.

% \paragraph{Step 3: different region labels imply proximity to an interface.}
% Let $r_t=\iota_\theta(t,A_t,x)$ and
% $r_t'=\iota_{\theta'}(t,A_t',x')$. By \eqref{eq:selection}, if $r_t'\ne r_t$,
% then at least one corresponding sign differs. For this index $j$,
% \[
%  |\ell_{\theta,j}(t,A_t,x)|
%  \le|\ell_{\theta',j}(t,A_t',x')-\ell_{\theta,j}(t,A_t,x)|.
% \]
% Using \eqref{eq:interface-regularity} and compactness of time and state,
% \[
%  |\ell_{\theta',j}(t,A_t',x')-\ell_{\theta,j}(t,A_t,x)|
%  \le C\|Z_t\|+C\delta(1+\|A_t\|)\le C\delta R.
% \]
% Outside a null set in $dt\,d\mathbb P$,
% \begin{equation}
%  \ind_{\{r_t'\ne r_t\}}
%  \le\sum_{j=1}^J
%        \ind_{\{|\ell_{\theta,j}(t,A_t,x)|\le C\delta R\}}.
%  \label{eq:switch-strip}
% \end{equation}

% \paragraph{Step 4: weighted random-width occupation.}
% For each $j$, partition according to $E_n=\{2^n\le R<2^{n+1}\}$,
% $n\ge0$. By Holder with exponents $4$ and $4/3$, \eqref{eq:strip}, and
% Markov's inequality,
% \begin{align*}
%  \E[R\,O_j(C\delta R)]
%  &\le\sum_{n\ge0}2^{n+1}
%        \E[O_j(C\delta2^{n+1})\ind_{E_n}]\\
%  &\le\sum_{n\ge0}2^{n+1}\|O_j(C\delta2^{n+1})\|_{L^4}
%        \mathbb P(R\ge2^n)^{3/4}\\
%  &\le C\delta\sum_{n\ge0}2^{2n}
%        \left(\frac{\E R^4}{2^{4n}}\right)^{3/4}
%  \le C\delta\sum_{n\ge0}2^{-n}\le C\delta.
% \end{align*}
% No independence between $R$ and the action path is used. By
% \eqref{eq:switch-strip} and $1+\|A_t\|\le R$,
% \begin{equation}
%  \E\int_0^T(1+\|A_t\|)\ind_{\{r_t'\ne r_t\}}dt\le C\delta.
%  \label{eq:weighted-switch}
% \end{equation}

% \paragraph{Step 5: compare parameter derivatives using the affine templates.}
% Define the $p\times d$ branch derivative by
% \[
%  (H_{\theta,S}(t,a,x))_{i,\cdot}
%  =K\left(\partial_{\theta_i}M^1_{\theta,S}t
%           +\partial_{\theta_i}M^2_{\theta,S}a
%           +\partial_{\theta_i}M^3_{\theta,S}x
%           +\partial_{\theta_i}b_{\theta,S}\right)^\top.
% \]
% The coefficient assumptions imply
% \begin{align}
%  \|H_{\theta,S}(t,a,x)\|_{\HS}&\le C(1+\|a\|),\\
%  \|H_{\theta',S}(t,a',x')-H_{\theta,S}(t,a,x)\|_{\HS}
%  &\le C\|a'-a\|+C(1+\|a\|)\delta.
% \end{align}
% Off switching hyperplanes, continuity of the tests makes the active label
% locally constant in $\theta$, so this branch derivative is the ordinary
% parameter derivative. Step 2 permits ignoring boundary times. Add and
% subtract $H_{\theta,r_t'}(t,A_t,x)$ to obtain
% \begin{align}
%  \|H_t'-H_t\|_{\HS}
%  &\le\|H_{\theta',r_t'}(t,A_t',x')-H_{\theta,r_t'}(t,A_t,x)\|_{\HS}\\
%  &\quad+\|H_{\theta,r_t'}(t,A_t,x)-H_{\theta,r_t}(t,A_t,x)\|_{\HS}\\
%  &\le C\|Z_t\|+C(1+\|A_t\|)\delta
%          +C(1+\|A_t\|)\1_{\{r_t'\ne r_t\}}.
% \end{align}
% The affine extension in the added term need not be active at the unprimed
% input; the template assumptions hold there regardless. Integrating,
% \eqref{eq:moment}, \eqref{eq:path-stability}, and
% \eqref{eq:weighted-switch} yield
% \begin{equation}
%  \int_0^T\E\|H_t'-H_t\|_{\HS}dt\le C\delta.
%  \label{eq:h-av}
% \end{equation}

% \paragraph{Step 6: include inverse volatility.}
% Write $\sigma_t=\sigma(t,A_t,x)$ and
% $\sigma_t'=\sigma(t,A_t',x')$. The inverse identity gives
% \[
%  \|\sigma_t'^{-1}-\sigma_t^{-1}\|_{\op}
%  \le\kappa^{-1}\|\sigma_t'-\sigma_t\|_{\op}
%  \le C(\|Z_t\|+\|x'-x\|).
% \]
% Consequently,
% \begin{align*}
%  \|G_t'-G_t\|_{\HS}
%  &\le\kappa^{-1/2}\|H_t'-H_t\|_{\HS}
%     +C(1+\|A_t\|)(\|Z_t\|+\|x'-x\|).
% \end{align*}
% Cauchy--Schwarz and the moment/stability estimates give
% \[
%  \E[(1+\|A_t\|)\|Z_t\|]
%  \le\|1+\|A_t\|\|_{L^2}\|Z_t\|_{L^2}\le C\delta.
% \]
% Integrating and using \eqref{eq:h-av} proves \eqref{eq:a7av}.
% The linear-growth bound on $H$ and bounded inverse volatility also ensure
% that $G$ is square integrable, so the displayed scores are well-defined.
% \end{proof}



% \begin{remark}[What is additional to Eq. (4.7)]
% For fixed architecture and consistently labelled activation templates,
% branch coefficients are polynomial in the weights, and the usual
% parameter perturbation estimate is available on bounded boxes. The new
% substantive hypothesis is Assumption C: a fixed switching description with
% uniformly controlled normalized coefficients and nonzero action normals.
% It excludes unrestricted degeneracy and unhandled state-only switching.
% Equation (4.7), the upper slope bounds, and bounded weights do not imply
% this condition. A sufficient way to verify the normalization is to start
% with raw affine tests whose coefficients are uniformly bounded and
% Lipschitz in $\theta$, and whose action-normal magnitudes are uniformly
% bounded below, then divide by those magnitudes. Tests with a fixed sign
% on the entire parameter/input domain can be dropped. Other degenerate
% tests require separate treatment or sharper cancellation-aware estimates.
% \end{remark}

% \begin{remark}[Why fourth-order occupation estimates appear]
% In the restricted shallow proof, the relevant gate jumps have bounded
% amplitudes and an unweighted $L^2$ occupation bound suffices. Here a branch
% parameter derivative can grow linearly in the action, so Step 4 controls
% $\E[R O_j(C\delta R)]$. The $L^4$ strip estimate closes this bound using
% only a fourth moment of $R$; no eighth moment of the initialization is
% needed. The constant depends on the finite switching description and is
% not asserted to have polynomial dependence on architecture size.
% \end{remark}

\newpage

\section{PL Condition}
identify subspace $U$

\begin{enumerate}
    % \item check $1+\|a\|$
    \item rectified flow
    \item check RELU compatiability
    % \item nextstep sample-based analysis
    % \item page 26 of the lecture notes
    \item check experiment to see variance reduction
    \item \textbf{Question: clipping}
\end{enumerate}
\bibliographystyle{plainnat}
\bibliography{references}
\end{document}
